% referee = single column, 1.75x line spacing for submission
% extra   = loads gji_extra.sty which provides natbib citation commands
\documentclass[extra]{gji}
%\documentclass[referee]{gji}

\usepackage{amsmath}
\usepackage{amssymb}
\usepackage{graphicx}
\graphicspath{{../../}}
\usepackage{booktabs}
\usepackage[colorlinks=true, allcolors=blue]{hyperref}
% gji_extra.sty embeds natbib but omits \citealt — add it here
\makeatletter
\DeclareRobustCommand\citealt
   {\begingroup\NAT@swafalse\def\NAT@ctype{0}\NAT@parfalse\NAT@citetp}
\makeatother
\usepackage{lineno}
\linenumbers

\DeclareMathOperator{\sgn}{sgn}
\newcommand{\RE}{\mathrm{Re}}
\newcommand{\IM}{\mathrm{Im}}
\newcommand{\infint}{\displaystyle\int_{-\infty}^{+\infty}}

% ====================================================================
\title[Navarrete \& Olive -- CSS inter-rifting model]
  {An analytical cyclo-steady-state model for inter-rifting deformation at magmatic rifts}

\author[Navarrete \& Olive]
  {Iv\'an Navarrete$^1$ \and Jean-Arthur Olive$^1$ \\
  $^1$Laboratoire de G\'eologie, \'Ecole Normale Sup\'erieure,
  PSL Research University, Paris, France
  }

\date{}
\pagerange{\pageref{firstpage}--\pageref{lastpage}}
\pubyear{\the\year{}}

% ====================================================================
\begin{document}
\label{firstpage}
\maketitle

% ====================================================================
\begin{summary}
Inter-rifting deformation (the strain that accumulates between discrete
magmatic episodes) remains the least-constrained phase of the rifting
cycle. Existing kinematic frameworks represent the rift as a continuously
opening dislocation in an elastic half-space, an assumption incompatible
with both the known rheological structure of magmatic rifts and far-field
geodetic constraints. We develop a fully analytical model for the surface
fields generated by a periodic sequence of instantaneous dike intrusions
in an elastic plate overlying a Newtonian viscous half-space. The
governing equation reduces, in the Fourier domain, to a fractional
(half-Laplacian) diffusion equation parameterized by a single
characteristic velocity $a = EH/[2\eta(1-\nu^2)]$, where $E$, $H$, and
$\nu$ are the Young's modulus, thickness, and Poisson's ratio of the
brittle layer, and $\eta$ is the asthenospheric viscosity. In the limit of
many repeated intrusions the system reaches a Cyclo-Steady State (CSS)
with a closed-form expression through the Digamma and log-Gamma functions,
conserving far-field plate divergence by construction. We extend the
model to a Maxwell visco-elastic asthenosphere via the
elastic--visco-elastic correspondence principle, recovering two physically
distinct relaxation regimes separated by the Maxwell relaxation time
$\tau_R$: near-field elastic rebound for $t \lesssim \tau_R$ and
far-field viscous diffusion for $t \gg \tau_R$. We fit the viscous and Maxwell CSS directly to two magmatic rift
systems, deriving $\tau_R$ from the same elastic modulus shared by plate
and half-space rather than treating it as an independent free parameter,
so the two rheologies are compared at equal parameter count. At Krafla
(Foulger et al., 1992), the viscous CSS yields
$\eta \approx 4\times10^{18}$~Pa~s ($H = 10$~km, $E = 40$~GPa), with
Maxwell rheology only marginally disfavoured ($\Delta\mathrm{AIC} \approx 1$);
at Asal \citep{vigny2007twenty}, the viscous CSS gives
$\eta \approx 2\times10^{19}$~Pa~s, a factor of five higher than at
Krafla, and is decisively favoured over its Maxwell counterpart
($\Delta\mathrm{AIC} \approx 30$) because the 25-yr record there samples
enough of the post-diking transient to rule out a spurious elastic
contribution.  The viscosity contrast between the two sites reflects the
anomalous sub-lithospheric heat flux beneath Iceland. The CSS framework
provides a mechanically self-consistent description of inter-rifting
deformation at any magmatic rift system, yielding viscosity estimates
from a single geodetic velocity profile without auxiliary rheological
assumptions.
\end{summary}

\begin{keywords}
continental rifting; inter-rifting deformation; fractional diffusion;
visco-elastic relaxation; GPS geodesy; Iceland
\end{keywords}

% ====================================================================
\section{Introduction}
\label{sec:intro}

Continental rifting drives the transition from stable cratonic lithosphere to active oceanic spreading through rifting cycles that span a broad range
of spatial and temporal scales \citep{corti2009continental,ebinger2013rifting}. These
cycles are conventionally partitioned into three phases: elastic strain
release through dike emplacement and normal faulting \citep{grandin2011seismicity,maccaferri2016time}; post-diking visco-elastic relaxation, which
decays on timescales of years to decades \citep{foulger1992post,nooner2009post}; and an inter-rifting phase during which the plates diverge at
a secular rate set by large-scale mantle dynamics \citep{saria2014present,stamps2021redefining}. Geodetic campaigns in Iceland \citep{heki1993plate,hofton1995horizontal} and the Afar \citep{vigny2007twenty,nooner2009post,grandin2011seismicity} have documented the first two phases in considerable detail. The
mechanisms governing the build-up of extensional strain between discrete
events, however, remain the least-constrained aspect of the rifting cycle.

The interseismic strain accumulation in rifts is poorly constrained \citep{ebinger2013rifting}.
Modelling the inter-rifting phase is conventionally done by transplanting kinematic frameworks developed for inter-seismic
earthquake cycle. 
The standard approach to this problem borrows the backslip model developed
for the interseismic earthquake cycle \citep{savage1978asthenosphere}. The rift
is assumed locked to depth $H$, below which it slips continuously,
generating the basal shear forces needed to reproduce observed velocity
profiles. Alternative formulations represent the rift as a
dislocation in an elastic half-space that opens at a depth-dependent rate
\citep{lafemina2005geodetic,konca2018slip}, or prescribe a basal shear
traction to mimic far-field plate forcing \citep{ali2014geodetic}. These models
share a fundamental physical inconsistency: they ignore the rheological
transition at the brittle--ductile boundary of magmatic rifts, where high
heat flow and partial melt weaken the lower crust and uppermost mantle
\citep{prodehl1997east,plasman2017lithospheric}. Geodetic and seismic
constraints place this transition at 5--15~km depth in Iceland and the
Afar \citep{stevens2021tectonic,weinstein2017fault}, implying visco-elastic
relaxation times incompatible with the purely elastic response assumed by
backslip models. A continuously opening dislocation is equally problematic:
dike intrusions are finite, discrete events, and the geodetic record shows
that each episode is followed by a localized transient decaying within a
few years \citep{foulger1992post,reiss2022shallow}.

Efforts to model the post-diking transient have confronted the same
boundary-condition obstacle from two directions. Analytical approaches,
exemplified by the diffusion model of \citet{foulger1992post}, reproduce the
broad deformation pattern following individual intrusions but require an
a priori specification of the viscous layer thickness as a free parameter
(Figure~\ref{fig:hofton_diffusion}). The same profile captures the transition
from post-rifting relaxation to steady inter-rifting divergence directly.
Twelve to seventeen years after the onset of rifting, the displacement field
already traces a broad, arctangent-like transition rather than a transient
confined to the rift. Both ends of the traverse settle onto distinct,
non-zero rates instead of returning to zero far from the rift axis.
Numerical visco-elastic simulations \citep{hofton1995horizontal,hofton1996post,nooner2009post} face an
equivalent difficulty: in any finite domain, velocities must vanish at the
lateral boundaries, a condition irreconcilable with the steady plate
divergence observed in global reference frames \citep{fernandes2004defining,altamimi2016itrf2014}. The result is a kinematic gap between the
transient post-diking signal and the secular strain rates of global plate
motion models \citep{kreemer2014geodetic,stamps2018virtual}.

\begin{figure}[htbp]
  \centering
  \includegraphics[width=\columnwidth]{fig/fig_hofton_krafla_diffusion.png}
  \caption{Figure recreated from \citet{hofton1996post}, their Figure 16:
    rift-perpendicular horizontal displacements at Krafla, 1987--1992, with
    $1\sigma$ error bars (black circles), and the \citet{foulger1992post}
    diffusion model $u(x,\tau) = b\,\mathrm{erfc}[x/(2\sqrt{\kappa\tau})]$,
    differenced between $\tau_i=12$ and $\tau_f=17$~yr after the onset of the
    1975--1984 rifting sequence, at the $\kappa=3.11$~m$^2$~s$^{-1}$,
    $b=1.27$~m best-fit parameters reported by \citeauthor{hofton1996post}
    (red curve; a constant reference-frame offset, not a free shape parameter
    of the model, is added to match the digitized profile). The systematic
    misfit at near- and far-field stations simultaneously ($\chi^2_r \approx
    10$) leads \citet{hofton1996post} to argue that a viscoelastic
    half-space, rather than simple diffusion, is needed to explain the
    post-rifting deformation field.}
  \label{fig:hofton_diffusion}
\end{figure}

The theoretical foundations laid for cyclic strike-slip fault deformation
\citep{nur1974postseismic,savage1978asthenosphere,spence1979magma,cohen1984crustal} suggest a different perspective. Inter-rifting
deformation is not a distinct physical process but the emergent
superposition of visco-elastic responses to an infinite history of discrete
dike intrusions. The apparent steadiness of inter-rifting velocities arises
naturally from the time-averaging of decaying transients, with no need for
a continuously creeping source.

We develop an analytical model for an elastic crust overlying a viscous
half-space \citep{hetland2006thesis} in which dike intrusions recur periodically at
interval $T$ and each open the rift by $b$, the dike opening. The governing
integro-differential equation reduces, after a Fourier transform, to a
fractional diffusion equation \citep{metzler2000random,samko1993integration} parameterized by a single characteristic velocity
$a = EH/[2\eta(1-\nu^2)]$. In the limit of many repeated intrusions, the
system reaches a Cyclo-Steady State (CSS) expressed exactly through the
Digamma and log-Gamma functions, which satisfies the far-field plate
divergence constraint by construction. We then extend the model to a
Maxwell visco-elastic asthenosphere via the elastic--visco-elastic
correspondence principle \citep{hetland2005postseismic,barbot2010unified}, recovering two physically distinct relaxation regimes separated by
the Maxwell timescale $\tau_R = \eta/\mu_m$. We apply both models to two
geodetic datasets: horizontal displacement profiles from the Krafla rift,
Iceland, recorded 8--11 years after the 1975--1984 rifting sequence
\citep{foulger1992post}, and a 25-year baseline time series from the
Asal rift, Djibouti \citep{vigny2007twenty}. Together, they constrain asthenospheric viscosity across
settings separated by a factor of five, and show that decadal post-rifting transient and secular plate divergence are the short- and long-term expression of a singly underlying process.
and secular plate divergence.

% ====================================================================
\section{Model}
\label{sec:model}

\subsection{Physical setup and governing equations}
\label{sec:setup}

The model represents the lithosphere as a coupled mechanical system: a thin
elastic plate of thickness $H$ (the seismogenic crust) overlying a
Newtonian viscous half-space of viscosity $\eta$ representing the asthenosphere
\citep[e.g.,][]{hetland2005postseismic,hetland2006thesis}.
The horizontal coordinate $x$ is perpendicular to the rift axis; the vertical
coordinate $z$ is oriented downward, with the free surface at $z = 0$ and the
crust--mantle interface at $z = H$. Under plane-strain conditions with
negligible vertical stress ($\sigma_{zz}\approx 0$), the depth-averaged
horizontal stress in the elastic layer is related to the horizontal surface
displacement $u(x,t)$ through Hooke's law:
\begin{equation}
\bar{\sigma}_{xx} = \frac{E}{1-\nu^2}\frac{\partial u}{\partial x},
\label{eq:hookean}
\end{equation}
where $E$ is Young's modulus and $\nu$ is Poisson's ratio. The factor
$(1-\nu^2)^{-1}$ accounts for the lateral confinement of plane strain.
Mechanical equilibrium of the plate requires the horizontal gradient of the
depth-averaged stress to be balanced by the shear traction exerted at its base
by the underlying viscous mantle:
\begin{equation}
H\frac{\partial\bar{\sigma}_{xx}}{\partial x} + \sigma_{xz}^{\text{BOT}} = 0.
\label{eq:equil}
\end{equation}

The rheological character of the asthenosphere enters through the constitutive
relation for the basal traction $\sigma_{xz}^{\text{BOT}}$. For a Newtonian
viscous half-space, the shear stress at $x$ depends on the surface velocity
at \emph{all} other points through the non-local Stokes kernel
\citep{savage1978asthenosphere}:
\begin{equation}
\sigma_{xz}^{\text{BOT}}(x,t) =
  \int_{-\infty}^{+\infty}\frac{2\eta}{\pi(x-x')^2}\,
  \frac{\partial u}{\partial t}(x',t)\,dx',
\label{eq:stokes}
\end{equation}
interpreted as Hadamard finite part. This non-local coupling reflects the
long-range character of viscous flow in the half-space: a velocity perturbation
at $x'$ decays as $(x-x')^{-2}$ rather than being confined to a finite
diffusion radius. In magmatically active rifts, where high heat flow and partial
melt weaken the lower crust \citep{holtzman2016questions,plasman2017lithospheric}, this
coupling operates across the entire plate boundary system.

\citet{foulger1992post} start from the same force balance but close it
differently. Rather than coupling the elastic layer to an infinite viscous
half-space through the non-local Stokes kernel (\ref{eq:stokes}), they place
a viscous channel of finite thickness $H_v$ over a rigid base. Its basal
traction is then treated as local, simple-shear drag,
$\sigma_{xz}^{\mathrm{BOT}} \sim (\eta/H_v)\,\partial u/\partial t$, in place
of the convolution in (\ref{eq:stokes}). Substituted into (\ref{eq:equil}),
this closure collapses the governing equation to an ordinary, Fickian
diffusion equation for the surface displacement, $\partial u/\partial t =
\kappa\,\partial^2 u/\partial x^2$. The diffusivity $\kappa \propto E H
H_v/\eta$ this introduces aggregates the elastic-layer properties with the
assumed viscous-channel thickness $H_v$, rather than constraining a
half-space viscosity directly. Reproducing a given deformation pattern
therefore requires an a priori choice of $H_v$ at each site -- the
limitation already noted in Section~\ref{sec:intro}. Section~\ref{sec:fourier}
derives what the non-local closure (\ref{eq:stokes}) yields in its place: an
infinite half-space, rather than a finite channel, replaces the ordinary
Laplacian of Foulger et al.'s diffusion equation with the half-order
fractional Laplacian $(-\Delta)^{1/2}$ of Equation~\ref{eq:fractional}. The
two models therefore differ in the order of the governing diffusion
equation itself, not only in its diffusivity.

Substituting
(\ref{eq:hookean}) into (\ref{eq:equil}) and using (\ref{eq:stokes}) yields
the governing integro-differential equation for $u$:
\begin{equation}
\frac{EH}{1-\nu^2}\frac{\partial^2 u}{\partial x^2}
+ \int_{-\infty}^{+\infty}\frac{2\eta}{\pi(x-x')^2}
  \frac{\partial u}{\partial t}(x',t)\,dx' = 0.
\label{eq:govPDE}
\end{equation}
Dike intrusions enter as initial conditions: each event opens the rift
instantaneously by $b$, offsetting the block on one side of the axis
relative to the other (Figure~\ref{fig:schematic}).
Section~\ref{sec:step} and Appendix~\ref{app:b2} show that this offset
subsequently redistributes into a symmetric $\pm b/2$ split about the
pre-intrusion position far from the rift, while the opening measured at the
intrusion site itself remains exactly $b$ at all times.

\begin{figure}[htbp]
  \centering
  \includegraphics[width=0.5\textwidth]{fig/EsquemaModelo.png}
  \caption{Model geometry.  The brittle elastic plate (thickness $H$,
    Young's modulus $E$, Poisson's ratio $\nu$) overlies a Newtonian viscous
    half-space (viscosity $\eta$).  Dike intrusions open the rift by
    $b$ at intervals $T$ (Appendix~\ref{app:b2}); the surface displacement
    $u(x,t)$ is positive away from the rift axis.  The parameter
    $a = EH/[2\eta(1-\nu^2)]$ sets the characteristic velocity that couples
    elastic and viscous deformation.}
  \label{fig:schematic}
\end{figure}

\subsection{Fourier reduction and the characteristic velocity}
\label{sec:fourier}

Equation (\ref{eq:govPDE}) is non-local in $x$ and not tractable in physical
space. A spatial Fourier transform resolves this, because the Stokes convolution
diagonalizes in wavenumber space. Defining the transform pair
\begin{equation*}
\hat{u}(k,t) = \int_{-\infty}^{+\infty}u(x,t)e^{-ikx}\,dx,
\qquad
u(x,t) = \frac{1}{2\pi}\int_{-\infty}^{+\infty}\hat{u}(k,t)e^{ikx}\,dk,
\end{equation*}
the transform of $\partial^2 u/\partial x^2$ is $-k^2\hat{u}$. The transform
of the Stokes kernel convolution is the classical result
\citep{savage1978asthenosphere}:
\begin{equation}
\mathcal{F}\!\left\{
  \int_{-\infty}^{+\infty}\frac{2\eta}{\pi(x-x')^2}
  \frac{\partial u}{\partial t}\,dx'
\right\} = -2\eta|k|\,\frac{\partial\hat{u}}{\partial t}.
\label{eq:stokes_FT}
\end{equation}
Applying the Fourier transform to (\ref{eq:govPDE}) and using
(\ref{eq:stokes_FT}):
\begin{equation}
-\frac{EH}{1-\nu^2}k^2\hat{u} - 2\eta|k|\frac{\partial\hat{u}}{\partial t} = 0.
\label{eq:fourier_pre}
\end{equation}
Rearranging and introducing the characteristic velocity
\begin{equation}
\boxed{a := \frac{EH}{2\eta(1-\nu^2)}},
\label{eq:a}
\end{equation}
equation (\ref{eq:fourier_pre}) simplifies to a decoupled first-order ODE for
each wavenumber $k$:
\begin{equation}
\frac{\partial\hat{u}}{\partial t} = -a\,|k|\,\hat{u}.
\label{eq:fourier}
\end{equation}
The parameter $a$ encodes the full elastic--viscous coupling of the system:
large $a$ (stiff crust, low viscosity) implies rapid redistribution of
post-diking strain over broad spatial scales; small $a$ implies slow, locally
confined relaxation. In physical space, the operator $|k|$ corresponds to the
fractional half-Laplacian $(-\Delta)^{1/2}$, so (\ref{eq:fourier}) is
equivalently written as the fractional diffusion equation
\citep{metzler2000random,samko1993integration,samko1993fractional}
\begin{equation}
\frac{\partial u}{\partial t} = -a\,(-\Delta)^{1/2}u.
\label{eq:fractional}
\end{equation}
This order is not a modelling choice. The fractional Laplacian admits a
representation as a Dirichlet-to-Neumann map for a harmonic extension
problem into an auxiliary half-space, weighted by $y^{1-2s}$; the weight
vanishes, and the extension becomes the ordinary Laplace equation, only at
$s=1/2$ \citep{caffarelli2007extension,daoud2022fractional}. That is the
same half-space potential problem generating the Stokes kernel itself, so
the half-order non-locality in (\ref{eq:fractional}) is forced by the
geometry of the elastic-plate--viscous-half-space coupling, not fitted to
the data. Solving the full two-dimensional Stokes flow beneath the plate,
governed by the biharmonic rather than the harmonic equation, still
recovers the same interface operator $-2\eta(-\Delta)^{1/2}$ exactly once
the interface stress is extracted.

This contrasts with classical Fickian diffusion: the half-order operator, a
direct consequence of the non-local Stokes kernel, produces Lorentzian rather
than Gaussian spatial profiles, with tails that decay as $1/x$ rather than
exponentially. The solution of (\ref{eq:fourier}) for arbitrary initial data is:
\begin{equation}
\hat{u}(k,t) = \hat{u}(k,0)\,e^{-a|k|t}.
\label{eq:fourier_sol}
\end{equation}

\subsection{Step response to a single dike intrusion}
\label{sec:step}

The simplest building block of the model is the deformation produced by a
single instantaneous dike intrusion at $t = 0$, which opens the rift by $b$.
Because the governing equation (\ref{eq:govPDE}) constrains $u$ only through
its gradient, the reference level from which this offset is measured is
arbitrary, and the initial condition can be written without loss of
generality as a Heaviside displacement:
\begin{equation}
u(x,0^+) = b\,\mathcal{H}(-x) = \begin{cases}
  0 \; , & x>0 \\
  b \; , & x<0
\end{cases}.
\label{eq:IC}
\end{equation}
Applying the Fourier transform to (\ref{eq:IC}) gives
$\hat{u}(k,0) = b\int_{-\infty}^{0}e^{-ikx}\,dx$,
which for real $k$ is treated in the distributional sense by adding a
convergence regulator $\varepsilon > 0$:
\begin{equation*}
\int_{-\infty}^{0}e^{-ikx}\,dx
= \lim_{\varepsilon\to 0}\int_{-\infty}^{0}e^{(-ik+\varepsilon)x}\,dx
= \lim_{\varepsilon\to 0}\frac{1}{\varepsilon - ik}.
\end{equation*}
By the Sokhotski--Plemelj theorem \citep{ablowitz2003complex},
\begin{equation}
\lim_{\varepsilon\to 0^+}\frac{1}{x \pm i\varepsilon} = \mathcal{P}\left(\frac{1}{x}\right) \mp i\pi\delta(x),
\label{eq:sokhotski}
\end{equation}
where $\mathcal{P}$ denotes the Cauchy principal value. Multiplying numerator
and denominator by $i$ recasts the regulator as $+i\varepsilon$,
\begin{equation*}
\frac{1}{\varepsilon - ik} = \frac{i}{i\varepsilon + k} = \frac{i}{k + i\varepsilon},
\end{equation*}
so that (\ref{eq:sokhotski}) applies directly with $x = k$ and the upper sign:
\begin{align}
\lim_{\varepsilon\to 0^+}\frac{1}{\varepsilon - ik}
&= \lim_{\varepsilon\to 0^+}\frac{i}{k + i\varepsilon} \notag\\
&= i\left[\mathcal{P}\left(\frac{1}{k}\right) - i\pi\delta(k)\right] \notag\\
&= \pi\delta(k) + \frac{i}{k}, \label{eq:sokhotski_applied}
\end{align}
where the principal-value symbol on $1/k$ is dropped from here on, as is
standard, since it is understood implicitly whenever $1/k$ appears next to a
$\delta(k)$ term. Hence the initial Fourier spectrum is:
\begin{equation}
\hat{u}(k,0) = b\pi\delta(k) + \frac{ib}{k}.
\label{eq:IC_FT}
\end{equation}
Substituting (\ref{eq:IC_FT}) into (\ref{eq:fourier_sol}) and noting that
$\delta(k)$ is active only where $e^{-a|k|t}=1$:
\begin{equation}
\hat{u}(k,t) = b\pi\delta(k) + \frac{ib}{k}\,e^{-a|k|t}.
\label{eq:evolved}
\end{equation}
The inverse Fourier transform splits into two contributions,
$u(x,t) = I_1 + I_2$. The delta-function term gives a constant background
$I_1 = b/2$; for $I_2$, the parity of $\sin(kx)/k$ yields:
\begin{equation}
I_2 = -\frac{b}{\pi}\int_0^{\infty}\frac{\sin(kx)}{k}\,e^{-akt}\,dk.
\label{eq:I2}
\end{equation}
This is a standard Dirichlet-type integral. Applying the complementary
arctangent identity $\arctan(z)+\arctan(1/z)=\pi/2$, the step displacement is:
\begin{equation}
\boxed{u^{\text{STEP}}(x,t) = \frac{b}{\pi}\arctan\!\left(\frac{at}{x}\right).}
\label{eq:step}
\end{equation}
Differentiating with respect to $t$, the step velocity is a Lorentzian
profile that broadens in time:
\begin{equation}
v^{\text{STEP}}(x,t) = \frac{b a}{\pi}\,\frac{x}{x^2 + a^2t^2}.
\label{eq:stepv}
\end{equation}
Both expressions satisfy the correct limits: $u^{\text{STEP}}\to b/2$ and
$v^{\text{STEP}}\to 0$ as $t\to\infty$, consistent with far-field plate
divergence. This limit reflects the profile's self-similarity:
$u^{\text{STEP}}$ depends on $x$ and $t$ only through the ratio $x/(at)$. At
any fixed $x$, that ratio vanishes as $t\to\infty$, so the profile relaxes to
the midpoint of its own jump, $b/2$, rather than to $b$. This is not a loss
of the imposed opening: Appendix~\ref{app:b2} shows that the discontinuity
at the intrusion site itself remains exactly $b$ for all $t>0$, so the
amplitude $b$ fitted in \S\ref{sec:krafla} and \S\ref{sec:asal} is directly
the physical dike-opening width, not half of it. The velocity peak
migrates outward as $at$, with amplitude decaying as $1/t$. On the connection between (\ref{eq:step}) and the
\citet{foulger1992post} diffusion model, see the fractional-Laplacian annex,
\S 3.2.

\subsection{Cyclic accumulation and the Cyclo-Steady State}
\label{sec:css}

\subsubsection{Finite sum over past intrusions}

A rift that has experienced $n = \lfloor t/T \rfloor$ dike intrusions at
regular intervals of $T$ years accumulates displacement by superposition of
individual step responses:
\begin{equation}
u(x,t) = \frac{b}{\pi}\sum_{j=1}^{n}
  \arctan\!\left(\frac{a(t-jT)}{x}\right).
\label{eq:finite_sum}
\end{equation}
Differentiating with respect to $t$, the corresponding velocity field is:
\begin{equation}
v(x,t) = \frac{b a}{\pi}\sum_{j=1}^{n}
  \frac{x}{x^2 + a^2(t-jT)^2}.
\label{eq:finite_sum_v}
\end{equation}

\subsubsection{The CSS limit and Digamma representation}

After many cycles the system enters a periodic state, the rift analog of the
cyclic steady state described for the interseismic earthquake cycle \citep{thatcher1984viscoelastic,hetland2006interseismic}. Introducing
$\tau\in[0,T)$ as the time elapsed since the most recent intrusion, the
Cyclo-Steady State (CSS) velocity is obtained by extending
(\ref{eq:finite_sum_v}) to an infinite history:
\begin{equation}
v_{\text{CSS}}(x,\tau) = \frac{b a}{\pi}
  \sum_{j=1}^{\infty}\frac{x}{x^2 + a^2(\tau-jT)^2}.
\label{eq:css_sum}
\end{equation}
Because this series converges absolutely, it admits a closed-form expression
via the Digamma function $\psi(z)$. Using the partial-fraction identity
\begin{equation*}
\frac{x}{x^2+a^2(\tau-jT)^2}
= \frac{1}{2i}\!\left[
    \frac{1}{a(\tau-jT)-ix} - \frac{1}{a(\tau-jT)+ix}
  \right],
\end{equation*}
factoring $-aT$ from each denominator and substituting into
(\ref{eq:css_sum}), the double sum takes the form of differences of Digamma
series. Using the representation $\psi(z+1) = -\gamma +
\sum_{n=1}^\infty(1/n - 1/(n+z))$, the Euler--Mascheroni constant $\gamma$ and
the $1/n$ terms cancel when the two conjugate series are differenced. With
$Z = \tau/T + ix/(aT)$ and the identity
$\psi(Z)-\psi(\bar{Z}) = 2i\,\IM\{\psi(Z)\}$,
the CSS velocity collapses to:
\begin{equation}
\boxed{v_{\text{CSS}}(x,\tau) =
  \frac{2b}{\pi T}\,\IM\!\left\{
    \psi\!\left(\frac{\tau}{T}+i\frac{x}{aT}\right)
  \right\}.}
\label{eq:vcss}
\end{equation}

\subsubsection{CSS displacement and far-field conservation}

Integrating (\ref{eq:vcss}) over $\tau$, and using the antiderivative
$d[\ln\Gamma(z)]/dz = \psi(z)$, the CSS displacement is:
\begin{equation}
\boxed{u_{\text{CSS}}(x,\tau) =
  \frac{2b}{\pi}\,\IM\!\left\{
    \ln\Gamma\!\left(\frac{\tau}{T}+i\frac{x}{aT}\right)
    -\ln\Gamma\!\left(i\frac{x}{aT}\right)
  \right\}.}
\label{eq:ucss}
\end{equation}
Over one complete cycle, the cumulative displacement at any fixed $x>0$ equals
the dike opening:
\begin{equation}
\int_0^T v_{\text{CSS}}(x,\tau)\,d\tau = b, \qquad \forall\,x>0.
\label{eq:conservation}
\end{equation}
This is not imposed as a boundary condition but follows analytically from the
Digamma representation: $\Gamma(1+iq) = iq\,\Gamma(iq)$ implies
$\IM\{\ln\Gamma(1+iq)-\ln\Gamma(iq)\} = \pi/2$, giving the integral $b$.
The CSS thus reconciles high-frequency magmatic pulses with secular divergence
without unphysical boundary conditions, satisfying the far-field constraint of
global plate motion models \citep{saria2014present,stamps2021redefining} by
construction.

The spatial width of the CSS deformation zone scales as $aT \propto EHT/\eta$.
Because the CSS depends on the rift parameters $(b, T)$ only through the
product $aT$, geodetic profiles of inter-rifting deformation directly constrain
the ratio $EH/\eta$, providing a new route to lithospheric characterization in
systems where $T$ can be independently estimated
(Figure~\ref{fig:css_behavior}a).  The finite-sum approximation
to the CSS converges to within the measurement noise of a typical geodetic
network after approximately 10 rifting cycles
(Figure~\ref{fig:css_behavior}b); for systems with fewer events on record, the
CSS should be replaced by the explicit sum in (\ref{eq:finite_sum}), and only
the step-response regime of Section~\ref{sec:step} applies.

\begin{figure}[htbp]
  \centering
  \includegraphics[width=\columnwidth]{fig/fig_css_behavior.png}
  \caption{Dimensionless CSS velocity and convergence from finite rifting
    histories.  All quantities are normalised by $b$ and $T$; distances
    by $aT$.
    \textbf{(a)}~CSS velocity $v_\mathrm{CSS}\cdot T/b$ versus
    $\hat{x}=x/(aT)$ at four inter-rifting phases $\hat{\tau}=\tau/T$.
    The velocity peak width is locked at $aT$ regardless of phase; all
    profiles converge to $v_\mathrm{plate} = b/T$ (dotted line) in the
    far field.
    \textbf{(b)}~Velocity at $\hat{\tau}=0.5$ for $N=1$, 3, and 10 dike
    events (solid curves) compared with the CSS limit (dashed).  The $N=1$
    profile does not recover the far-field plate velocity; convergence to the
    CSS requires $\sim\!10$ cycles.}
  \label{fig:css_behavior}
\end{figure}

\subsection{Extension to a Maxwell visco-elastic asthenosphere}
\label{sec:maxwell_model}

\subsubsection{Constitutive model and the correspondence principle}

The Maxwell relaxation time $\tau_R := \eta/\mu_m$ (where $\mu_m$ is
the mantle shear modulus) separates two regimes: on timescales
$t \ll \tau_R$ the asthenosphere responds elastically; on timescales
$t \gg \tau_R$ it flows viscously \citep{burgmann2008rheology}.
We model this by replacing the Newtonian half-space with a Maxwell body: a
Hookean spring and a Newtonian dashpot in series. The 1D constitutive relation
for deviatoric stress $\tau_{ij}$ and strain $\varepsilon$ is:
\begin{equation}
\dot{\varepsilon} = \frac{\dot{\tau}}{2\mu_m} + \frac{\tau}{2\eta}.
\label{eq:maxwell_constitutive}
\end{equation}
To adapt the Stokes kernel (\ref{eq:stokes}) we employ the elastic--visco-elastic
correspondence principle \citep{hetland2005postseismic,barbot2010unified}:
in the Laplace domain, static moduli are replaced by their complex,
frequency-dependent counterparts. Applying the Laplace transform
$\tilde{f}(s)=\mathcal{L}\{f(t)\}$ to (\ref{eq:maxwell_constitutive}) with zero
initial conditions:
\begin{equation}
s\tilde{\varepsilon} = \tilde{\tau}\!\left(\frac{s}{2\mu_m}+\frac{1}{2\eta}\right).
\label{eq:laplace_ve}
\end{equation}
Solving for stress gives $\tilde{\tau} = 2\tilde{\mu}(s)\tilde{\varepsilon}$,
where the effective frequency-dependent shear modulus of the Maxwell body is:
\begin{equation}
\tilde{\mu}(s) = \frac{\mu_m\eta\,s}{\mu_m + \eta s}
               = \mu_m\frac{\tau_R s}{1+\tau_R s}.
\label{eq:maxwell_mu}
\end{equation}

\subsubsection{The visco-elastic Stokes kernel}

For a purely elastic half-space the basal traction relates to surface
displacement through a kernel proportional to $\mu_m$:
\begin{equation}
\sigma_{xz}^{\text{BOT, el}}(x) =
  \int_{-\infty}^{+\infty}\frac{2\mu_m}{\pi(x-x')^2}\,u(x')\,dx'.
\label{eq:elastic_kernel}
\end{equation}
Applying the correspondence principle (replacing $\mu_m$ with $\tilde{\mu}(s)$
in the Laplace domain) and factoring out $\eta = \mu_m\tau_R$:
\begin{equation}
(1+\tau_R s)\,\tilde{\sigma}_{xz}^{\text{BOT}}(x,s) =
  \int_{-\infty}^{+\infty}\frac{2\eta}{\pi(x-x')^2}
  \left[s\tilde{u}(x',s)\right]dx'.
\label{eq:ve_kernel_laplace}
\end{equation}
Inverting to the time domain (multiplication by $s$ in the Laplace domain
corresponds to $\partial/\partial t$) yields the visco-elastic Stokes kernel:
\begin{equation}
\sigma_{xz}^{\text{BOT}} + \tau_R\frac{\partial\sigma_{xz}^{\text{BOT}}}{\partial t}
= \int_{-\infty}^{+\infty}\frac{2\eta}{\pi(x-x')^2}
  \frac{\partial u(x',t)}{\partial t}\,dx'.
\label{eq:ve_kernel}
\end{equation}
In the limit $\tau_R\to 0$, (\ref{eq:ve_kernel}) reduces to the purely viscous
Stokes kernel (\ref{eq:stokes}).

\subsubsection{Governing equation in Fourier space}

The equilibrium equation (\ref{eq:equil}) expresses the basal traction in terms
of the elastic crust displacement:
\begin{equation}
\sigma_{xz}^{\text{BOT}} = -\frac{EH}{1-\nu^2}\frac{\partial^2 u}{\partial x^2}.
\label{eq:sigma_from_equil}
\end{equation}
Substituting (\ref{eq:sigma_from_equil}) and its time derivative into
(\ref{eq:ve_kernel}), then applying the spatial Fourier transform
using (\ref{eq:stokes_FT}):
\begin{equation}
\frac{EH}{1-\nu^2}k^2\hat{u}
+ \tau_R\frac{EH}{1-\nu^2}k^2\frac{\partial\hat{u}}{\partial t}
= -2\eta|k|\frac{\partial\hat{u}}{\partial t}.
\label{eq:maxwell_fourier_pre}
\end{equation}
Dividing by $2\eta$, using (\ref{eq:a}), and collecting
$\partial\hat{u}/\partial t$ terms gives the visco-elastic analogue of
(\ref{eq:fourier}):
\begin{equation}
\frac{\partial\hat{u}}{\partial t}
= -\left(\frac{a|k|}{1+\tau_R a|k|}\right)\hat{u}.
\label{eq:maxwell_fourier}
\end{equation}
The denominator $1+\tau_R a|k|$ acts as a low-pass filter on the relaxation
rate, slowing the decay of high-wavenumber (near-field) features relative to
the purely viscous case.

\subsubsection{Step response and asymptotic behavior}

Equation (\ref{eq:maxwell_fourier}) is separable in $k$, yielding a decaying
exponential for each wavenumber:
\begin{equation}
\hat{u}(k,t) = \hat{u}(k,0)\exp\!\left(-\frac{a|k|}{1+\tau_R a|k|}\,t\right).
\label{eq:maxwell_fourier_sol}
\end{equation}
Applying the Heaviside initial condition (\ref{eq:IC_FT}):
\begin{equation}
\hat{u}(k,t) = b\pi\delta(k)
  + \frac{ib}{k}\exp\!\left(-\frac{a|k|}{1+\tau_R a|k|}\,t\right).
\label{eq:maxwell_IC_evolved}
\end{equation}
Taking the inverse Fourier transform and following the same parity argument as
in Section~\ref{sec:step}, the Maxwell step displacement requires numerical
integration:
\begin{equation}
u^{\text{MAX}}(x,t) = \frac{b}{2}
  - \frac{b}{\pi}\int_0^{+\infty}
    \frac{\sin(kx)}{k}\exp\!\left(-\frac{akt}{1+\tau_R ak}\right)dk,
\label{eq:maxwell_step}
\end{equation}
since the rational exponent prevents the closed-form Dirichlet evaluation of
Section~\ref{sec:step}. The surface velocity follows by differentiating
(\ref{eq:maxwell_step}) with respect to $t$:
\begin{equation}
v^{\text{MAX}}(x,t) = \frac{b a}{\pi}\int_0^{+\infty}
  \frac{\sin(kx)}{1+\tau_R ak}
  \exp\!\left(-\frac{akt}{1+\tau_R ak}\right)dk.
\label{eq:maxwell_vel}
\end{equation}
Both integrals are evaluated numerically using adaptive Gauss-Kronrod
quadrature \citep{laurie1997calculation}, with integration breakpoints placed at the
successive zeros of $\sin(kx)$ (i.e.\ $k = n\pi/x$, $n = 1,2,\ldots$)
to partition the oscillatory integrand into well-behaved sub-intervals.
The upper truncation wavenumber is set adaptively to ensure the
exponential factor is negligible at the integration cutoff.

The exponent in (\ref{eq:maxwell_step}) reveals two physically distinct regimes.
In the long-wavelength limit ($k\to 0$ or $t\to\infty$), $1+\tau_R ak \approx 1$
and the exponential reduces to $\exp(-akt)$, recovering the purely viscous
Dirichlet integral of Section~\ref{sec:step}. The elastic memory of the
asthenosphere fades at large distances and late times, and the model converges
to the Newtonian viscous limit. In the short-wavelength limit ($k\to\infty$ or
$t\to 0$), $1+\tau_R ak \approx \tau_R ak$ and the exponent saturates:
\begin{equation}
\lim_{k\to\infty}\left(-\frac{akt}{1+\tau_R ak}\right) = -\frac{t}{\tau_R}.
\label{eq:maxwell_saturation}
\end{equation}
At early times $t\ll\tau_R$, the exponential approaches unity and no near-field
deformation occurs outside the rift interface. For $t\sim\tau_R$, near-field
deformation decays exponentially at the fixed rate $1/\tau_R$, representing
pure Maxwell relaxation of the mantle independent of the elastic crust.
Together, (\ref{eq:maxwell_vel}) and (\ref{eq:maxwell_saturation}) demonstrate
that the Maxwell model introduces a short-time, near-field stiffening relative
to the viscous model: both converge to the same Cyclo-Steady State at
$t\gg\tau_R$ (Section~\ref{sec:maxwell_css}).

\subsection{Maxwell Cyclo-Steady State}
\label{sec:maxwell_css}

\subsubsection{Infinite sum and Fourier periodicity condition}

The Maxwell CSS is the superposition of step-response velocities from all
preceding intrusions, by the same argument as in Section~\ref{sec:css}:
\begin{equation}
  v_{\text{CSS}}^{\text{MAX}}(x,\tau)
  = \sum_{j=0}^{\infty} v^{\text{STEP,MAX}}(x,\,\tau + jT).
\end{equation}
Writing $v^{\text{STEP,MAX}}$ in Fourier space
(Section~\ref{sec:maxwell_model}) as
$\frac{2b a}{\pi}\int_0^\infty \sin(kx)\,e^{-\alpha(k)(\tau+jT)}\,dk$,
with the effective Maxwell decay rate
\begin{equation}
  \alpha(k) := \frac{ak}{1+\tau_R ak},
  \label{eq:alpha_k}
\end{equation}
each past intrusion contributes a factor $e^{-\alpha(k)jT}$.  The series
converges for all $k>0$ and the geometric sum
$\sum_{j=0}^{\infty}e^{-\alpha(k)jT} = (1-e^{-\alpha(k)T})^{-1}$ can be
performed inside the Fourier integral.  The same result follows directly
from the requirement that each Fourier mode satisfy the periodicity
condition $\hat{u}(k,\tau+T) = \hat{u}(k,\tau) + ib/k$, which
admits the unique solution
\begin{equation}
  \hat{u}(k,\tau) = \frac{ib/k}{1-e^{-\alpha(k)T}}\,e^{-\alpha(k)\tau},
  \qquad \tau\in[0,T).
  \label{eq:maxwell_css_fourier}
\end{equation}

\subsubsection{Physical-space representation, conservation, and limits}

Inverting (\ref{eq:maxwell_css_fourier}) gives the Maxwell CSS velocity
and displacement:
\begin{align}
  v_{\text{CSS}}^{\text{MAX}}(x,\tau)
  &= \frac{2b a}{\pi}\int_0^{\infty}
     \frac{\sin(kx)}{1+\tau_R ak}\,
     \frac{e^{-\alpha(k)\tau}}{1-e^{-\alpha(k)T}}\,dk,
  \label{eq:maxwell_css_v} \\
  u_{\text{CSS}}^{\text{MAX}}(x,\tau)
  &= \frac{2b}{\pi}\int_0^{\infty}
     \frac{\sin(kx)}{k}\,
     \frac{1-e^{-\alpha(k)\tau}}{1-e^{-\alpha(k)T}}\,dk.
  \label{eq:maxwell_css_u}
\end{align}
Integrating the velocity over one full period, using
$\int_0^T e^{-\alpha\tau}/(1-e^{-\alpha T})\,d\tau = 1/\alpha$ and the
algebraic identity $a/[(1+\tau_R ak)\,\alpha(k)] = 1/k$, gives
\begin{equation}
  \int_0^T v_{\text{CSS}}^{\text{MAX}}(x,\tau)\,d\tau
  = \frac{2b}{\pi}\int_0^{\infty}
    \frac{\sin(kx)}{k}\,dk = b,
\end{equation}
where the last step is the Dirichlet integral.  The full dike opening
is recovered every cycle, independent of $\tau_R$.  In the limit
$\tau_R\to 0$, $\alpha(k)\to ak$ and $(1+\tau_R ak)^{-1}\to 1$, and
(\ref{eq:maxwell_css_v}) reduces to the viscous CSS Fourier integral whose
closed form is the Digamma series of Section~\ref{sec:css}.  For
$\tau_R > 0$, the saturating denominator $1+\tau_R ak$ prevents this
reduction and (\ref{eq:maxwell_css_v}) is evaluated numerically.

\subsubsection{Constraining $\tau_R$ without an extra free parameter}
\label{sec:tauR_shared}

Section~\ref{sec:results} fits the Maxwell CSS to real data alongside its
purely viscous counterpart.  Leaving $\tau_R$ free alongside $(a,\,b)$ would
let the extra parameter absorb noise rather than test whether Maxwell
rheology genuinely improves the fit, since a model with an additional free
parameter is guaranteed to reach at least as low a $\chi^2$ regardless of
whether that parameter is physically warranted.  We instead derive $\tau_R$
from $a$ itself, so that the Maxwell CSS carries exactly as many free
parameters as the viscous CSS it is compared against.

The half-space's Maxwell shear modulus $\mu_m$ (Equation~\ref{eq:maxwell_mu})
has so far been left general.  We adopt the simplest closure available
without additional data: the half-space shares the same shear modulus as the
elastic plate it underlies,
\begin{equation}
\mu_m = \mu_{\text{plate}} := \frac{E}{2(1+\nu)}.
\label{eq:mu_shared}
\end{equation}
Combined with $\eta = EH/[2a(1-\nu^2)]$, the inversion of the characteristic
velocity (\ref{eq:a}), this gives
\begin{equation}
\boxed{\tau_R = \frac{\eta}{\mu_m} = \frac{H}{a(1-\nu)}},
\label{eq:tau_R_shared}
\end{equation}
a relaxation time set entirely by the already-fitted characteristic velocity
$a$ and the canonical geometric parameters $H$, $\nu$
(Section~\ref{sec:disc_rheology}), with no independent estimate of $\mu_m$
required.  This is a simplifying assumption, not a physical necessity -- the
mantle shear modulus need not equal the crust's -- but it is the natural
default in the absence of independent constraints, and it keeps the
viscous-versus-Maxwell comparison in Section~\ref{sec:results} an honest,
equal-parameter test of functional form rather than one biased toward the
model with more free parameters.

% ====================================================================
\section{Results: Applications to magmatic rift systems}
\label{sec:results}

Post-diking geodetic records from two contrasting rift environments test the
models developed in Section~\ref{sec:model} and constrain sub-lithospheric
viscosity across distinct tectonic settings.  We apply the viscous step
response (Section~\ref{sec:step}) to a rift-perpendicular differenced
displacement profile from the 1975--1984 Krafla rifting sequence in NE
Iceland \citep{foulger1992post} and to a 25-year baseline time series
combined with a rift-perpendicular velocity profile from the 1978
Asal--Ghoubbet episode in the Afar Triangle \citep{vigny2007twenty}.
Both datasets fall within the single-event step-response regime: at
8--11~yr post-rifting, Krafla has experienced less than one recurrence
interval, and the Asal record ends before a subsequent major rifting
episode.  Together, the two sites span a factor-of-five contrast in
best-fit characteristic velocity $a$ and probe the rheological contrast
between a plume-fed mid-ocean rift and an early-stage continental rift.

\subsection{Dimensionless form and model behavior}
\label{sec:dimless}

The viscous and Maxwell step responses each admit a natural non-dimensionalization
that reveals the governing control parameters and delineates the physically
distinct relaxation regimes.

\subsubsection{Viscous case: a universal single-parameter curve}

Introducing the dimensionless distance $\hat{x} = x/(at)$ (the spatial
coordinate measured in units of the instantaneous diffusion front $at$),
the step displacement (\ref{eq:step}) collapses to
\begin{equation}
\frac{u^{\text{STEP}}}{b}
= \frac{1}{\pi}\arctan\!\left(\frac{1}{\hat{x}}\right).
\label{eq:step_dimless}
\end{equation}
Equation~(\ref{eq:step_dimless}) is a universal curve: every viscous step
response, regardless of the values of $E$, $H$, $\eta$, or $b$, maps
onto this single function of $\hat{x}$.  Near the rift axis or at late times
($\hat{x}\ll 1$), the displacement saturates at $u/b\to 1/2$; far from
the rift or at early times ($\hat{x}\gg 1$), $u/b\approx
1/(\pi\hat{x})\to 0$.  The corresponding velocity profile is an equally
compact Lorentzian,
\begin{equation}
\frac{at\,v^{\text{STEP}}}{b}
= \frac{1}{\pi}\,\frac{\hat{x}}{1+\hat{x}^2},
\label{eq:stepv_dimless}
\end{equation}
with a peak located precisely at the diffusion front ($\hat{x}=1$,
i.e.\ $x = at$), an amplitude that decays as $b/(at)$, and a width
that broadens linearly in time.  The single parameter $a$ thus encodes the
entire spatio-temporal evolution of the post-diking transient.

The surface profile is only the boundary trace of a flow that fills the
underlying half-space. Depth is reintroduced through the same harmonic
extension used to construct the fractional Laplacian in
Section~\ref{sec:fourier}: at depth $z$, the horizontal and vertical
interior velocity components are the Poisson-kernel convolution of the
surface velocity $V(x)$,
\begin{equation}
v_x(x,z) = \frac{2}{\pi}\int_{-\pi/2}^{\pi/2} \sin^2\theta\, V(x-z\tan\theta)\,d\theta,
\qquad
v_z(x,z) = \frac{2}{\pi}\int_{-\pi/2}^{\pi/2} \sin\theta\cos\theta\, V(x-z\tan\theta)\,d\theta,
\label{eq:halfspace_kernel}
\end{equation}
with $V$ set to the viscous step velocity $v^{\text{STEP}}(x,t)$ or the
viscous CSS velocity $v_{\text{CSS}}(x,\tau)$ as appropriate. Figure~\ref{fig:viscous_halfspace}
shows both, for a representative set of rheological parameters, at six
illustrative phases $\tau$ spanning early (elastic-like) to late (diffused)
states: the surface profile (top row of each block) alongside the interior
flow it drives (bottom row, arrows coloured and scaled by speed). Both
blocks are normalised by the same reference velocity $V$, so arrow colour
is directly comparable between STEP and CSS. In both models, near-field
flow is steep, shallow, and directed away from the rift at early $\tau$;
as $\tau$ grows, the same total flux spreads over a broadening and
deepening lobe, exactly the widening self-similar front of
Equations~(\ref{eq:step_dimless})--(\ref{eq:stepv_dimless}). The STEP
panels show the true decay of the transient, with the peak surface
velocity falling steadily from the shortest to the longest $\tau$ shown.
The CSS interior field (bottom block) never empties out this way: it
saturates on the far-field plate rate $V$ by construction
(Section~\ref{sec:css}), while the STEP field (top block) keeps thinning
toward zero.

\begin{figure*}[htbp]
  \centering
  \includegraphics[width=\textwidth]{fig/fig_halfspace_step_cycle.png}
  \caption{Surface velocity profile (top row of each block) and interior
    velocity field (bottom row; arrows coloured by speed magnitude, via
    Equation~\ref{eq:halfspace_kernel}) for the viscous STEP response (top
    two rows) and the viscous CSS (bottom two rows), at six illustrative
    recurrence times $\tau$. Both rows are normalised by the same reference
    velocity $V$ (the plate rate implied by the CSS cycle period), so arrow
    colour is directly comparable between STEP and CSS; arrow length is
    scaled independently within each block for visual clarity. The dashed
    vertical line in each surface-profile panel marks the diffusion front
    $L=a\tau$. The single grid point directly above the dike is omitted
    from the quiver because the near-source field there is disproportionately
    large relative to the rest of the domain.}
  \label{fig:viscous_halfspace}
\end{figure*}

\subsubsection{Maxwell case: two dimensionless groups and three regimes}

The Maxwell step solution introduces a second dimensionless group.  For
positive wavenumbers $k>0$, writing the Fourier-space exponent
(\ref{eq:maxwell_fourier_sol}) as
\begin{equation}
-\frac{akt}{1+\tau_R ak}
= -\frac{\hat{t}\,\kappa}{1+\kappa},
\label{eq:maxwell_dimless_exp}
\end{equation}
where $\hat{t} = t/\tau_R$ is the Maxwell number and $\kappa = \tau_R ak$ is
the wavenumber scaled by the Maxwell length
\begin{equation}
\ell_M := a\tau_R = \frac{EH}{2\mu_m(1-\nu^2)},
\label{eq:maxwell_length}
\end{equation}
three physically distinct relaxation regimes emerge.
(i)~\emph{Elastic rebound} ($\kappa\gg 1$, i.e.\ $x\ll\ell_M$): the
exponent in (\ref{eq:maxwell_dimless_exp}) saturates at $-\hat{t}$,
independent of wavenumber, so near-field deformation decays at the uniform
rate $1/\tau_R$, pure Maxwell relaxation of the mantle spring,
independent of the elastic crust.
(ii)~\emph{Viscous diffusion} ($\hat{t}\gg 1$ or $x\gg\ell_M$, $\kappa\ll 1$):
the denominator $1+\kappa\approx 1$ and the Maxwell solution converges
pointwise to the purely viscous step response of
Section~\ref{sec:step}.
(iii)~\emph{Crossover} ($\hat{t}\sim 1$): both timescales are simultaneously
active; near-field deformation lags behind the viscous prediction (a
short-time stiffening) before relaxing toward the common Cyclo-Steady
State at $t\gg\tau_R$.
The Maxwell length $\ell_M = EH/[2\mu_m(1-\nu^2)]$ is a \emph{purely
elastic} quantity, independent of viscosity: it measures the ratio of
lithospheric to mantle elastic stiffness times the plate thickness, and
sets the spatial boundary between the two asymptotic regimes.

Regime (i) admits a closed-form limit that sharpens this picture. As
$\kappa\to\infty$ the exponent in (\ref{eq:maxwell_dimless_exp}) saturates
at $-\hat{t}$ independently of $k$, so the high-wavenumber tail of the
Maxwell velocity integral (\ref{eq:maxwell_vel}) reduces to the same
Dirichlet kernel $\sin(kx)/k$ used to derive the viscous step response
(\ref{eq:I2}), now weighted by the $x$-independent factor
$e^{-\hat{t}}/(\tau_R a)$.  Evaluating that Dirichlet integral gives the
near-field plateau
\begin{equation}
v^{\text{MAX}}(x\to0^+,\,t) \;\longrightarrow\;
\frac{b}{2\tau_R}\,e^{-t/\tau_R}
\;=\; \frac{b}{t}\,\frac{\hat{t}}{2}\,e^{-\hat{t}},
\label{eq:maxwell_plateau}
\end{equation}
an $x$-independent value rather than a smoothly rising-then-decaying
Lorentzian peak: at fixed $t$, near-axis velocity saturates at a level set
purely by how much elastic relaxation has occurred, not by distance from
the rift.  Equation~(\ref{eq:maxwell_plateau}) is maximised at
$\hat{t}=1$, i.e.\ when the post-diking observation time equals the
Maxwell relaxation time.  For $\hat{t}\gg1$ (fast-relaxing mantle,
$\tau_R\ll t$) the elastic phase has already decayed away by the time of
observation and the plateau vanishes; for $\hat{t}\ll1$ (slow-relaxing,
near-elastic mantle) the spring has barely begun to release and the
plateau is again small.  Because the plateau is independent of $x$ while
the viscous Lorentzian (\ref{eq:stepv_dimless}) rises from zero at the
rift axis, the two curves necessarily cross at some $x_c(\tau_R)<\ell_M$:
Maxwell velocity exceeds the viscous prediction for $x<x_c$ and falls
below it for $x_c<x<\ell_M$.  The crossover moves toward the rift axis as
$\tau_R$ increases, so the near-field enhancement of
(\ref{eq:maxwell_plateau}) is confined to a shrinking neighbourhood of the
rift for an increasingly elastic mantle.

Figure~\ref{fig:maxwell_regimes} illustrates this behaviour for the step
response.  Both the snapshot time $t$ and the Maxwell length $\ell_M
= a\tau_R$ are anchored to the shared-modulus-closure relaxation time of
Equation~(\ref{eq:tau_R_shared}), evaluated at the same (Natron best-fit)
$a$ used for every curve in the figure -- not an additional assumption,
but the one $\tau_R$ already singled out by that closure for this $a$.
Seven Maxwell relaxation times spanning two orders of magnitude around
this closure value (from one tenth to ten times it) are plotted against
$x/\ell_M$ on a logarithmic axis to resolve the near-field structure.
Each Maxwell curve flattens into the plateau of
(\ref{eq:maxwell_plateau}) well inside $x<\ell_M$, crosses the purely
viscous Lorentzian at the marked point $x_c$, and converges onto it in
the far field.  The plateau height is non-monotonic in $\tau_R$, peaking
when $\tau_R$ matches the closure value and decaying for both smaller and
larger $\tau_R$, consistent with (\ref{eq:maxwell_plateau}); the crossover
moves from $x_c\approx0.92\,\ell_M$ at the smallest $\tau_R$ shown to
$x_c\approx0.14\,\ell_M$ at the largest.
Figure~\ref{fig:maxwell_displacement} traces the cumulative displacement
at $x = \ell_M$ through the crossover: Maxwell models lag the viscous
prediction for $t \lesssim \tau_R$ before converging to the common
asymptote $b/2$.  Together, the two figures define an observational
criterion.  Post-diking geodetic records collected within a few $\tau_R$
of the rifting event, and within the shrinking near-axis zone
$x<x_c(\tau_R)$, sit in the elastic-rebound regime, where the purely
viscous step response underestimates near-field velocity; beyond $x_c$
but still inside $\ell_M$, the same viscous model overestimates it.  Once
$t \gg \tau_R$ the elastic spring has fully relaxed and the simpler
viscous model is an accurate substitute everywhere.

\subsubsection{Cyclo-Steady State: the deformation-width scaling}

The CSS introduces a different natural length scale.  In the viscous CSS
(\ref{eq:vcss}), the Digamma argument is $\tau/T + ix/(aT)$, so the
deformation width scales as $aT$ rather than $at$.  Defining
$\hat{x}_{\text{CSS}} = x/(aT)$ and $\hat{\tau} = \tau/T \in [0,1)$,
the CSS velocity collapses to
\begin{equation}
\frac{T\,v_{\text{CSS}}}{b}
= \frac{2}{\pi}\,\IM\!\left\{
    \psi\!\left(\hat{\tau} + i\hat{x}_{\text{CSS}}\right)
  \right\},
\label{eq:vcss_dimless}
\end{equation}
a universal function of $(\hat{\tau},\,\hat{x}_{\text{CSS}})$ alone.
Unlike the step response, the CSS velocity profile does not broaden with
time: its width is locked at $aT$ by the periodicity of the forcing.
For the Maxwell CSS (\ref{eq:maxwell_css_v}), the ratio $\tau_R/T$ enters
as a third dimensionless group.  When $\tau_R/T \ll 1$ the Maxwell CSS
converges pointwise to (\ref{eq:vcss_dimless}); when $\tau_R/T \sim 1$,
near-field stiffening persists into the steady state, narrowing the
velocity peak relative to the purely viscous prediction
(Figure~\ref{fig:maxwell_regimes}).

\begin{figure}[htbp]
  \centering
  \includegraphics[width=\columnwidth]{fig/fig_maxwell_velocity.png}
  \caption{Maxwell visco-elastic near-field velocity, at the snapshot time
    and for the Maxwell length $\ell_M = a\tau_R$ set by the
    shared-modulus-closure relaxation time of
    Equation~(\ref{eq:tau_R_shared}) at the Natron best-fit $a$ -- not an
    independently assumed value, but the one already singled out by that
    closure for this $a$.  Distances are normalised by $\ell_M$, velocities
    by $b/\tau_R$ evaluated at that same closure value.  Grey shading marks
    the elastic-rebound zone
    ($x < \ell_M$).  Seven values of $\tau_R$, spanning one tenth to ten
    times the closure value, plus the purely viscous limit (black), which
    peaks at $x = \ell_M$ (dashed line).  Each Maxwell curve flattens into
    the $x$-independent near-field plateau of
    Equation~(\ref{eq:maxwell_plateau}) (dotted, same colour), crosses the
    viscous Lorentzian at the marked point $x_c(\tau_R)$, and converges onto
    it in the far field.  The plateau height is largest when $\tau_R$
    matches the closure value and decreases for $\tau_R$ away from it in
    either direction, while $x_c$ decreases monotonically as $\tau_R$
    increases.}
  \label{fig:maxwell_regimes}
\end{figure}

\begin{figure}[htbp]
  \centering
  \includegraphics[width=\columnwidth]{fig/fig_maxwell_displacement.png}
  \caption{Maxwell visco-elastic cumulative displacement.  Displacement is
    normalised by $b$.  Grey shading marks the elastic-rebound zone ($t$
    below the shared-modulus-closure relaxation time of
    Equation~\ref{eq:tau_R_shared} at the Natron best-fit $a$).  Cumulative
    displacement at $x = \ell_M$ for the viscous model (black) and two
    Maxwell cases ($\tau_R$ set to that closure value, orange; to five
    times it, dashed red).  Orange shading marks the displacement deficit
    relative to viscous.  Both Maxwell models converge to the common
    asymptote $b/2$ (dotted) after $t \gg \tau_R$.}
  \label{fig:maxwell_displacement}
\end{figure}

Figure~\ref{fig:maxwell_halfspace} extends this comparison to the interior
flow field, mirroring Figure~\ref{fig:viscous_halfspace} but for the
Maxwell rheology: velocities are normalised by $\tau_R/b$ rather than
$1/V$, and the fixed Maxwell length $\ell_M$
(Equation~\ref{eq:maxwell_length}) replaces the growing diffusion front
$a\tau$ as the vertical guide in every panel. The Maxwell STEP interior
field (top two rows) inherits the near-field plateau of
Equation~(\ref{eq:maxwell_plateau}): peak surface velocity is largest at
the shortest $\tau$ shown and falls with increasing $\tau$, and the
underlying flow lobe shrinks and thins in step, the same qualitative decay
as the viscous STEP field but paced by $\tau_R$ rather than by the
diffusion front. The Maxwell CSS interior field (bottom two rows) traces
the same evolution within a single recurrence interval: immediately after
an intrusion ($\tau$ small) the profile is sharply peaked, the
elastic-rebound signature of Equation~(\ref{eq:maxwell_plateau}); as
$\tau$ grows toward $T$ the peak broadens and flattens, the spring
relaxing toward the smoother, viscous-like profile it would reach were no
further dike to intervene. This within-cycle relaxation is the spatial
counterpart of the near-field stiffening described above
($\tau_R/T\sim1$): the Maxwell CSS never fully forgets the elastic
response of its most recent intrusion, in contrast with the viscous CSS,
which saturates cleanly on the far-field plate rate throughout the cycle
(Figure~\ref{fig:viscous_halfspace}).

\begin{figure*}[htbp]
  \centering
  \includegraphics[width=\textwidth]{fig/fig_halfspace_step_cycle_maxwell.png}
  \caption{Surface velocity profile (top row of each block) and interior
    velocity field (bottom row; arrows coloured by speed magnitude, via
    Equation~\ref{eq:halfspace_kernel}) for the Maxwell STEP response (top
    two rows) and the Maxwell CSS (bottom two rows), at six illustrative
    recurrence times $\tau$. Both rows are normalised by $\tau_R/b$ (the
    Maxwell relaxation time), on the same colour scale so arrow colour is
    directly comparable between STEP and CSS; arrow length is scaled
    independently within each block for visual clarity. The dashed
    vertical line in each surface-profile panel marks the fixed Maxwell
    length $\ell_M = a\tau_R$ (Equation~\ref{eq:maxwell_length}), rather
    than the viscous $L=a\tau$, since $\ell_M$ does not depend on $\tau$.
    The single grid point directly above the dike is omitted from the
    quiver because the near-source field there is disproportionately
    large relative to the rest of the domain.}
  \label{fig:maxwell_halfspace}
\end{figure*}

At each site, model parameters are estimated by minimizing the chi-squared misfit
\begin{equation}
\chi^2 = \sum_{i=1}^{N}
  \left(\frac{u_i^{\mathrm{obs}} - u_i^{\mathrm{pred}}}{\sigma_i}\right)^2,
\label{eq:chi2}
\end{equation}
where $u_i^{\mathrm{obs}}$ and $u_i^{\mathrm{pred}}$ are the observed and
predicted displacement (or baseline change) at data point $i$, $\sigma_i$ is
the $1\sigma$ measurement uncertainty, and $N$ is the number of observations.
Goodness of fit is quantified through the reduced statistic
$\chi^2_r = \chi^2/(N-p)$, where $p$ is the number of free parameters;
$\chi^2_r \approx 1$ indicates that the model reproduces the data within the
measurement noise.

\subsection{Krafla rift, Iceland}
\label{sec:krafla}

The 1984 Krafla rifting episode in NE Iceland provides one of the
best-documented post-diking geodetic records available.  \citet{foulger1992post} measured horizontal surface displacements at 20 stations across
the rift zone at two epochs 8 and 11 years after the 1975--1984 rifting
sequence.  We fit the differenced displacement field
$\Delta u = u(\tau_f) - u(\tau_i)$, with $\tau_i = 8$~yr and
$\tau_f = 11$~yr, using the CSS forward models of
Sections~\ref{sec:css} and~\ref{sec:maxwell_css}: the viscous CSS, with
free parameters $(a,\,b)$ at fixed recurrence interval $T$ and with $T$
itself free; and its Maxwell counterpart, with $\tau_R$ derived from $a$
through the shared-modulus closure (\ref{eq:tau_R_shared}) rather than
fit independently, so that each Maxwell variant carries exactly as many
free parameters as its viscous counterpart.  Parameters are estimated by
minimizing equation~(\ref{eq:chi2}) over a Latin Hypercube sample in
log-uniform parameter space, followed by bounded Nelder-Mead refinement.

Leaving $T$ free returns $T \approx 3800$--$4300$~yr for the viscous and
Maxwell CSS respectively -- one to two orders of magnitude longer than
the $\sim250$~yr gap between the documented 1724--1729 and 1975--1984
eruptive crises at Krafla (Section~\ref{sec:disc_rheology}).  This is not
a measurement of the true recurrence interval: at 8--11~yr post-rifting
the observations span only a tiny fraction of any plausible $T$, so the
fit is free to push $T$ arbitrarily large, toward the single-event limit
in which the CSS reduces to the step response of
Section~\ref{sec:step}.  We therefore also fix $T = 200$~yr, consistent
with the documented recurrence history, as the physically anchored
comparison.

All four CSS variants fit the observed displacement profile within the
measurement noise (Figure~\ref{fig:krafla}, Table~\ref{tab:krafla_css}).
The viscous CSS gives $a = 5.17\times10^{-5}$~m~s$^{-1}$,
$b = 0.59$~m ($\chi^2_r = 0.72$) at $T = 200$~yr, and
$a = 4.87\times10^{-5}$~m~s$^{-1}$, $b = 0.59$~m,
$T = 4289$~yr ($\chi^2_r = 0.76$) with $T$ free -- both close to the
best-fit characteristic velocity obtained from the step response alone,
as expected since $\tau_i,\tau_f \ll T$ in every variant.  The Maxwell
CSS gives $a = 1.42\times10^{-4}$~m~s$^{-1}$, $b = 0.50$~m,
$\tau_R = 2.99$~yr (derived) at $T = 200$~yr ($\chi^2_r = 0.78$), and
$a = 1.34\times10^{-4}$~m~s$^{-1}$, $b = 0.46$~m,
$T = 3789$~yr, $\tau_R = 3.15$~yr (derived) with $T$ free
($\chi^2_r = 0.83$).  Because each Maxwell fit is compared against a
viscous fit with the same number of free parameters, the AIC comparison
reduces to a direct $\chi^2$ contrast: $\Delta\mathrm{AIC} = 1.1$
($T = 200$~yr) and $1.2$ ($T$ free), in both cases favouring the viscous
CSS but far short of the $|\Delta\mathrm{AIC}| > 2$ threshold for
positive evidence.  As at the single-event step response
(Section~\ref{sec:step}), the 8--11~yr Krafla record does not require a
Maxwell asthenosphere: whatever elastic contribution the mantle
half-space might make relaxes faster than the shortest interval sampled
here.  Using the viscous CSS best-fit $a$ ($T = 200$~yr) in~(\ref{eq:a})
yields $\eta = EH/[2a(1-\nu^2)] \approx 4.1\times10^{18}$~Pa~s for
canonical Icelandic lithospheric parameters ($H = 10$~km, $E = 40$~GPa,
$\nu = 0.25$), consistent with post-glacial rebound estimates
\citep{sigmundsson2010geodetic} and below the upper bound
$\eta < 10^{19}$~Pa~s placed by \citet{sigmundsson1991postglacial} for a 10-km
elastic layer at Krafla.

\begin{table}
\centering
\caption{CSS model comparison for the Krafla displacement profile
  ($N = 20$).  MaxCSS carries $\tau_R$ derived from $a$
  (\ref{eq:tau_R_shared}), not as an independent free parameter, so $k$
  matches the corresponding viscous CSS row.}
\label{tab:krafla_css}
\begin{tabular}{lccc}
\hline
Model & $k$ & $\chi^2_r$ & $\Delta$AIC vs.\ viscous \\
\hline
CSS ($T$ free)     & 3 & 0.76 & --   \\
CSS ($T=200$~yr)   & 2 & 0.72 & --   \\
MaxCSS ($T$ free)  & 3 & 0.83 & $+1.2$ \\
MaxCSS ($T=200$~yr)& 2 & 0.78 & $+1.1$ \\
\hline
\end{tabular}
\end{table}

\begin{figure}[htbp]
  \centering
  \includegraphics[width=\columnwidth]{fig/fig08_fits_css_comparison.png}
  \caption{Best-fit CSS differenced displacement profiles for the Krafla
    post-diking dataset of \citet{foulger1992post}.  Observed
    rift-perpendicular displacements between $\tau_i = 8$~yr and
    $\tau_f = 11$~yr post-rifting are shown as filled circles with
    $1\sigma$ error bars (cm).  Curves show the viscous CSS ($T$ free,
    purple; $T=200$~yr, dashed orange) and Maxwell CSS ($T$ free, navy;
    $T=200$~yr, dashed brown), with $\tau_R$ derived from $a$ via
    Equation~(\ref{eq:tau_R_shared}); see Table~\ref{tab:krafla_css} and
    text for parameter values and goodness-of-fit statistics.}
  \label{fig:krafla}
\end{figure}

Figure~\ref{fig:krafla_chi2} shows the $\chi^2$ confidence region for the
viscous CSS $(a,b)$ at $T=200$~yr, the best-fitting model at this site.
The $1\sigma$ contour is broad and elongated along an
anti-correlated $a$--$b$ trade-off: a larger characteristic velocity can
be compensated by a smaller dike opening without degrading the fit,
reflecting the modest station coverage of the single differenced epoch
pair.

\begin{figure}[htbp]
  \centering
  \includegraphics[width=\columnwidth]{fig/fig_chi2_landscape_krafla_css.png}
  \caption{$\chi^2/N$ confidence region for the viscous CSS fit to the
    Krafla displacement profile ($T=200$~yr, Table~\ref{tab:krafla_css}).
    Contours mark the 1$\sigma$, 2$\sigma$, and 3$\sigma$ confidence
    levels; the star marks the best fit $(a^*,b^*)$.}
  \label{fig:krafla_chi2}
\end{figure}

\subsection{Asal rift, Djibouti}
\label{sec:asal}

The November 1978 Asal--Ghoubbet rifting episode in the Afar Triangle opened
approximately 2~m of rift-perpendicular extension over several days \citep{ledain1979seismotectonic,ruegg1984seismotectonic}, launching one of the longest continuously
monitored post-diking transient records outside Iceland.  \citet{vigny2007twenty}
tracked the EP00--DF00 baseline (a rift-perpendicular geodetic tie crossing
the Asal rift) from 1979 to 2003, recording $\approx490$~mm of cumulative
lengthening, with rapid early growth decelerating sharply after the mid-1980s.
We digitize all 33 data points from Figure~3 of \citet{vigny2007twenty}, combining
trilateration campaigns (1979--1987) and GPS surveys (1991--2003), and fit the
same four CSS variants as at Krafla (Section~\ref{sec:krafla}): viscous CSS
with $T$ free or fixed, and its Maxwell counterpart with $\tau_R$ derived from
$a$ through (\ref{eq:tau_R_shared}).  The
predicted baseline change is evaluated as $\Delta L = 2u(x_0,\tau)$ at
$x_0 = 2000$~m, half the EP00--DF00 rift-perpendicular separation, with the
rifting onset fixed at $t_0 = 1978.85$~yr \citep{vigny2007twenty}.  Because the
\citeauthor{vigny2007twenty} time series is referenced to the first measurement epoch
($\tau_{\mathrm{ref}} \approx 1.1$~yr after the rifting event) rather than to the
rifting onset, we subtract from each model curve its own prediction at
$\tau_{\mathrm{ref}}$, anchoring $\Delta L = 0$ at the reference epoch without
introducing a free parameter.  No documented rifting episode preceded 1978 at
Asal, so we adopt $T = 500$~yr as the physically anchored fixed-$T$ comparison,
following the same reasoning as at Krafla.

All four CSS variants reproduce the overall shape of the Asal time series
(Figure~\ref{fig:asal_ts}, Table~\ref{tab:asal_css}), but the viscous CSS fits
markedly better than its Maxwell counterpart at equal parameter count.  The
viscous CSS gives $a = 9.47\times10^{-6}$~m~s$^{-1}$, $b = 0.354$~m
($\chi^2_r = 1.54$) at $T = 500$~yr, and $a = 9.70\times10^{-6}$~m~s$^{-1}$,
$b = 0.354$~m, $T = 2085$~yr ($\chi^2_r = 1.54$) with $T$ free -- as at
Krafla, the free-$T$ fit drifts far above any plausible recurrence interval,
confirming that the 25-yr Asal record cannot constrain $T$ independently
of the single-event limit.  The Maxwell CSS gives $a = 8.08\times10^{-5}$~m~s$^{-1}$,
$b = 0.316$~m, $\tau_R = 5.23$~yr (derived) at $T = 500$~yr
($\chi^2_r = 2.61$), and $a = 7.63\times10^{-5}$~m~s$^{-1}$, $b = 0.325$~m,
$T = 4066$~yr, $\tau_R = 5.54$~yr (derived) with $T$ free
($\chi^2_r = 2.48$).  Unlike at Krafla, the AIC comparison here is decisive:
$\Delta\mathrm{AIC} = 33.4$ ($T = 500$~yr) and $28.1$ ($T$ free), both
strongly favouring the viscous CSS ($|\Delta\mathrm{AIC}| \gg 6$).  The
25-yr Asal baseline, unlike the single 8--11~yr Krafla epoch pair, spans
enough of the post-diking transient that a spurious elastic contribution
from the derived $\tau_R \approx 5$--$6$~yr measurably worsens the fit at
late times, where the data are already well described by pure viscous
relaxation.

Independent field observations, which favoured the viscous single-event
response over the \citet{foulger1992post} diffusion model in earlier
analyses of this same record, do not resolve as cleanly for the CSS.  The
viscous CSS dike opening $b = 0.354$~m predicts, via the same $\Delta
L = 2u(x_0,\tau)$ rift-crossing relation used in the forward model
(\S\ref{sec:asal}), an EP00--DF00 baseline swing $2b = 0.71$~m per
cycle, and the Maxwell CSS gives $2b = 0.63$~m; both fall further
below the $\approx1.9$~m measured across the rift by trilateration
immediately after the 1978 event \citep{ruegg1979geodetic} than a
single-event fit unconstrained by the rest of the 25-yr record.  The same
pattern appears in the early deformation rate: the viscous CSS predicts
$\approx33$~mm~yr$^{-1}$ at $\tau_{\mathrm{ref}} \approx 1.1$~yr and the
Maxwell CSS $\approx39$~mm~yr$^{-1}$, both below the $\approx53$~mm~yr$^{-1}$
measured by trilateration shortly after the event \citep{ruegg1987deformation}.
This is a consequence of the CSS parameterisation itself: tying the
plate rate to $v_{\text{plate}} = b/T$ links $b$ to the late-time behaviour
of the whole 25-yr record, not only to its rapid early rise, and pulls the
best-fit $b$ below the value an early-epoch-only fit would return.  The
viscous best-fit $a = 9.47\times10^{-6}$~m~s$^{-1}$ nonetheless implies a
sub-lithospheric viscosity of $\eta \approx 2.3\times10^{19}$~Pa~s for Afar
lithospheric parameters ($H = 10$~km, $E = 40$~GPa, $\nu = 0.25$), in the
range of estimates from other post-diking transient analyses in the East
African Rift \citep{nooner2009post}.

\citet{vigny2007twenty} also measured rift-perpendicular GPS velocities along a
transect crossing the rift at a mean epoch of $\approx19$~yr after the 1978
event.  We forward-predict the velocity profile from the four CSS
time-series best fits, placing the model centre at $x = +1$~km northeast
of the topographic rift axis to coincide with the Petit Rift
\citep{vigny2007twenty}.  All four variants reproduce the near-rift
velocity gradient (Figure~\ref{fig:asal_profile}).  The dominant misfit is
the $\approx16$~mm~yr$^{-1}$ Somalia--Danakil background plate motion
\citep{saria2014present,stamps2021redefining}, which produces a pronounced
SW--NE velocity asymmetry absent from our post-rift transient models; we
fit a uniform background velocity ($V_{\mathrm{plate}} \approx 8$~mm~yr$^{-1}$
in every variant) to each profile before comparing shapes.  The viscous
CSS again fits better than its Maxwell counterpart, and slightly better
than the single-event step response reported in earlier analyses of this
profile: $\chi^2_r = 15.7$ ($T$ free) and $15.9$ ($T=500$~yr) for the
viscous CSS, versus $22.6$ ($T$ free) and $22.2$ ($T=500$~yr) for the
Maxwell CSS.  None of the four variants reaches $\chi^2_r \sim 1$; as
before, the elevated residuals reflect the asymmetric far-field
Somalia--Danakil signal that a single symmetric transient and a uniform
velocity offset cannot fully capture, not a failure specific to either
rheology.

\begin{table}
\centering
\caption{CSS model comparison for the Asal EP00--DF00 time series
  ($N = 33$).  MaxCSS carries $\tau_R$ derived from $a$
  (\ref{eq:tau_R_shared}), matching $k$ to the corresponding viscous CSS
  row.}
\label{tab:asal_css}
\begin{tabular}{lccc}
\hline
Model & $k$ & $\chi^2_r$ & $\Delta$AIC vs.\ viscous \\
\hline
CSS ($T$ free)     & 3 & 1.54 & --    \\
CSS ($T=500$~yr)   & 2 & 1.54 & --    \\
MaxCSS ($T$ free)  & 3 & 2.48 & $+28.1$ \\
MaxCSS ($T=500$~yr)& 2 & 2.61 & $+33.4$ \\
\hline
\end{tabular}
\end{table}

\begin{figure}[htbp]
  \centering
  \includegraphics[width=\columnwidth]{fig/fig_asal_timeseries_css_comparison.png}
  \caption{Best-fit CSS forward models for the Asal post-diking time series of
    \citet{vigny2007twenty}.  The EP00--DF00 baseline change (mm) is shown as
    filled circles with $1\sigma$ error bars.  Curves show the viscous CSS
    ($T$ free, purple; $T=500$~yr, dashed orange) and Maxwell CSS ($T$ free,
    navy; $T=500$~yr, dashed brown), with $\tau_R$ derived from $a$ via
    Equation~(\ref{eq:tau_R_shared}).  All curves are referenced to the first
    measurement epoch ($\tau_{\mathrm{ref}} \approx 1.1$~yr); see
    Table~\ref{tab:asal_css} and text for parameter values and
    goodness-of-fit statistics.}
  \label{fig:asal_ts}
\end{figure}

\begin{figure}[htbp]
  \centering
  \includegraphics[width=\columnwidth]{fig/fig_asal_velocity_profile_css.png}
  \caption{Forward-predicted rift-perpendicular velocity profiles at
    $t\approx19$~yr post-rifting, compared with GPS observations of Vigny
    et al.\ (2007) (Somalia reference frame, mm~yr$^{-1}$).  Curves show
    the viscous CSS ($T$ free, purple; $T=500$~yr, dashed orange) and
    Maxwell CSS ($T$ free, navy; $T=500$~yr, dashed brown), with
    parameters fixed at the time-series best-fit values
    (Figure~\ref{fig:asal_ts}) and a uniform plate-motion offset
    ($V_{\mathrm{plate}} \approx 8$~mm~yr$^{-1}$) added to each curve.
    Shading marks the 1978 activated zone ($\pm3$~km, light) and the Petit
    Rift ($\pm1$~km, dark); the dashed vertical line marks the model
    centre at $x = +1$~km.}
  \label{fig:asal_profile}
\end{figure}

Figure~\ref{fig:asal_chi2} shows the $\chi^2$ confidence region for the
viscous CSS $(a,b)$ at $T=500$~yr, the best-fitting model at this site.
Unlike Krafla, the 25-yr time series confines $(a,b)$ to a narrow,
strongly anti-correlated ellipse: the longer baseline resolves the
characteristic velocity and dike opening far more tightly, which is
precisely why the Maxwell CSS is decisively rejected here
(Table~\ref{tab:asal_css}) but not at Krafla.

\begin{figure}[htbp]
  \centering
  \includegraphics[width=\columnwidth]{fig/fig_chi2_landscape_asal_css.png}
  \caption{$\chi^2/N$ confidence region for the viscous CSS fit to the
    Asal time series ($T=500$~yr, Table~\ref{tab:asal_css}).  Contours
    mark the 1$\sigma$, 2$\sigma$, and 3$\sigma$ confidence levels; the
    star marks the best fit $(a^*,b^*)$.  Note the much narrower range
    spanned by both axes relative to Figure~\ref{fig:krafla_chi2}.}
  \label{fig:asal_chi2}
\end{figure}

% ====================================================================
\section{Discussion}
\label{sec:discussion}

% ── §4.1 ─────────────────────────────────────────────────────────────────────
\subsection{A boundary-condition-free description of inter-rifting deformation}
\label{sec:disc_bcs}

The CSS framework closes the kinematic gap identified in
Section~\ref{sec:intro} without introducing a locking depth, a prescribed
domain width, or an artificially imposed far-field velocity.  Backslip models
of inter-rifting deformation \citep{savage1978asthenosphere} and their dislocation
equivalents \citep{lafemina2005geodetic,konca2018slip} reproduce observed
velocity gradients only at the cost of a free geometric parameter (the
locking depth or creep-zone geometry) that carries no independent mechanical
constraint in magmatic rifts where the rheological transition is gradational
rather than sharp.  Numerical visco-elastic simulations face an equivalent
obstacle: in any finite domain, lateral boundary conditions must either force
far-field velocities to zero, in contradiction with global plate motion models
\citep{altamimi2016itrf2014}, or prescribe the divergence rate that the model is
itself meant to predict \citep{hofton1995horizontal,hofton1996post}.

The CSS requires neither.  The deformation zone half-width
$aT = EHT/[2\eta(1-\nu^2)]$ emerges directly from the balance between elastic
restoring stresses in the plate and viscous drag at its base, parameterized
entirely by the characteristic velocity $a$ and the recurrence interval $T$.
Far-field plate divergence is satisfied by construction through the conservation
property of the Digamma representation (equation~\ref{eq:conservation}): no
boundary condition is imposed, and no artificial velocity offset is needed to
reconcile model predictions with far-field geodesy.  Both free parameters are
physically grounded: the dike opening $b$ is measurable independently
from field geology and geodetic surveys \citep{ledain1979seismotectonic,ruegg1984seismotectonic}, and $a$ directly constrains the ratio $EH/\eta$ without requiring an
independent estimate of viscous layer thickness.  This structural simplicity
distinguishes the CSS from existing kinematic and semi-analytical approaches
to the inter-rifting problem.

% ── §4.2 ─────────────────────────────────────────────────────────────────────
\subsection{Rheological constraints and deformation zone scaling}
\label{sec:disc_rheology}

The single composite parameter $a = EH/[2\eta(1-\nu^2)]$ collapses the
elastic--viscous trade-off onto a one-dimensional family of solutions in
$(E,\,H,\,\eta)$ space.  Figure~\ref{fig:rheology_1sigma} makes this
trade-off explicit: for fixed $E$, the iso-$a$ contour traces a band across
$(H,\,\eta)$ space.  Dark shading propagates the 1$\sigma$ uncertainty on $a$
from the $\chi^2$ fit at fixed $E = 40$~GPa; light shading additionally
propagates $E$ over $[20,\,80]$~GPa (combined 1$\sigma$ envelope).  The iso-$a$
curves are nearly parallel in this representation, so the dominant uncertainty
is in $\eta$ rather than in $H$.  Independent seismological constraints on the
seismogenic layer thickness ($H \approx 5$--$15$~km in both Iceland and the
Afar; \citealt{stevens2021tectonic,plasman2017lithospheric,weinstein2017fault} break
the remaining ambiguity and allow us to translate the geodetically constrained
$a$ into a viscosity estimate for each site
(Figure~\ref{fig:rheology_1sigma}).

Adopting $H = 10$~km and $E = 40$~GPa, the Krafla viscous-CSS best-fit
$a = 5.17\times10^{-5}$~m~s$^{-1}$ ($T=200$~yr, Section~\ref{sec:krafla})
implies $\eta \approx 4.1\times10^{18}$~Pa~s, essentially the same estimate
the single-event step response would give, since $\tau_i,\tau_f \ll T$ places
the fit in the regime where the two forward models coincide.  Two
independent sources of uncertainty affect this estimate.  The 1$\sigma$
confidence interval on $a$ from the $\chi^2$ fit (Section~\ref{sec:krafla})
yields $\eta \in [2.6,\,7.4]\times10^{18}$~Pa~s for fixed $E = 40$~GPa, while
propagating $E$ over $[20,\,80]$~GPa at the best-fit $a$ gives
$\eta \in [2.1,\,8.3]\times10^{18}$~Pa~s.  Because $\eta = EH/[2a(1-\nu^2)]$,
the two uncertainties are additive in logarithmic space; combining them in
quadrature ($\sigma_{\log_{10} E} = 0.30$, $\sigma_{\log_{10} a} = 0.23$)
yields a combined 1$\sigma$ range of
$\eta \in [1.7,\,10.0]\times10^{18}$~Pa~s
(Figure~\ref{fig:rheology_1sigma}, left).  This range lies entirely below or at
the upper bound $\eta < 10^{19}$~Pa~s placed by \citet{sigmundsson1991postglacial} for a 10-km
elastic layer at Krafla, and is consistent with post-glacial rebound estimates
across Iceland ($\eta \sim 10^{18}$--$10^{19}$~Pa~s; \citealt{sigmundsson2010geodetic})
and with forward visco-elastic modelling of the same post-diking transient
\citep{hofton1996post}.  The low viscosity beneath Iceland reflects
anomalously high temperatures in the sub-lithospheric mantle maintained by the
Iceland plume: partial melt in the lower crust and uppermost mantle reduces the
effective viscosity by one to two orders of magnitude relative to standard
continental lithosphere \citep{holtzman2016questions}.

The Asal viscous-CSS best-fit $a = 9.47\times10^{-6}$~m~s$^{-1}$
($T=500$~yr, Section~\ref{sec:asal}) implies $\eta \approx 2.3\times10^{19}$~Pa~s
at $H = 10$~km and $E = 40$~GPa, with an $E$-propagated range of
$\eta \in [1.1,\,4.5]\times10^{19}$~Pa~s, approximately five times higher
than the Krafla estimate.  This is higher
than viscosity estimates for the Dabbahu segment of the same rift system
($\eta \approx 5\times10^{18}$~Pa~s; \citealt{nooner2009post};
$\eta \sim 10^{18}$~Pa~s, $H = 21$~km; Hamling et al., 2014 [CITE NEEDED]), reflecting
either genuine segment-to-segment variation in sub-lithospheric temperature
or contributions from shallow magmatic sources that the purely viscous CSS
cannot distinguish from mantle relaxation.  Smittarello et al.\\ (2016) [CITE NEEDED]
argue, from 38~years of geodetic observations of the Asal--Ghoubbet rift,
that magma supply rather than viscoelastic relaxation controls the
post-diking velocity field at this site, a degeneracy our single-parameter
model cannot resolve without independent constraints on the magmatic budget.
The factor-of-five viscosity contrast between Krafla and Asal nevertheless
remains physically coherent: both are magmatically active rifts, but
Iceland's plume-driven heat flux maintains substantially higher
sub-lithospheric temperatures, accelerating viscous relaxation relative to
the East African setting.

Figure~\ref{fig:css_width} translates the best-fit $a$ values into
predictions for the CSS deformation zone width $aT$, which controls how broadly
inter-rifting velocities are distributed across a geodetic network.  For a
Krafla recurrence interval of $T \approx 200$~yr, consistent with the
$\sim250$~yr gap between the 1724--1729 and 1975--1984 eruptive crises, the CSS
width reaches $aT \approx 326$~km, already exceeding the 230~km span of the
Foulger et al.\ (1992) network (Figure~\ref{fig:css_width}), consistent with the broad
velocity gradients observed across Iceland \citep{lafemina2005geodetic}.  For Asal the recurrence interval is
poorly constrained; adopting $T \approx 500$~yr gives $aT \approx 149$~km,
nearly ten times the 17~km aperture of the \citet{vigny2007twenty} GPS profile
(Figure~\ref{fig:css_width}).  In both cases the CSS width vastly exceeds the
network aperture, which is precisely why inter-rifting velocity fields in
Iceland and the Afar have been interpreted as steady plate divergence rather than
as the superposition of slowly decaying visco-elastic transients.  This is
also why the free-$T$ fits of Sections~\ref{sec:krafla}
and~\ref{sec:asal} cannot recover $T$ itself: once $aT$ exceeds the
network aperture by an order of magnitude, the data constrain only $a$,
and $\chi^2$ is essentially flat with respect to $T$ beyond that point.

\begin{figure}[htbp]
  \centering
  \begin{tabular}{@{}cc@{}}
    \includegraphics[width=4.0cm]{fig/fig_rheology_1sigma.png} &
    \includegraphics[width=4.0cm]{fig/fig_rheology_1sigma_asal.png} \\
  \end{tabular}
  \caption{Rheological trade-off for Krafla (left, red) and Asal (right, blue),
    using the single-event step-response $\chi^2$ landscape; the viscous-CSS
    best-fit $a$ (Sections~\ref{sec:krafla}, \ref{sec:asal}) is statistically
    indistinguishable ($<1\%$ different) at both sites, so the same iso-$a$
    curve and uncertainty envelope apply.  Each panel shows the iso-$a$
    contour in the $(H,\,\eta)$ plane at the
    best-fit value.  Dark shading propagates the 1$\sigma$ uncertainty on $a$
    from the $\chi^2$ fit at fixed $E = 40$~GPa; light shading additionally
    propagates $E$ over $[20,\,80]$~GPa (combined 1$\sigma$ envelope).  The
    $H = 10$~km dashed line marks the reference elastic thickness; markers show
    published viscosity estimates (Krafla: open circle, \citealt{sigmundsson1991postglacial}; open
    square, \citealt{hofton1996post}; open circle with error bar, \citealt{wright2012geophysical}; Asal: filled diamond, \citealt{nooner2009post} for the Dabbahu segment of the
    same rift system; orange box, \citealt{cattin2005numerical} — near-axis effective
    viscosity $\eta = 0.5$--$1\times10^{18}$~Pa~s from a 2D thermo-mechanical
    FEM model of the first 6--8~yr post-rifting at Asal, with the seismogenic
    layer thickness $H = 3$--$12$~km from the same study).  The Cattin et al.\
    box lies well below the iso-$a^*$ curve, consistent with the local
    near-axis viscosity being lower than the laterally averaged bulk estimate
    from the 25-yr geodetic record.}
  \label{fig:rheology_1sigma}
\end{figure}

\begin{figure}[htbp]
  \centering
  \includegraphics[width=\columnwidth]{fig/fig_css_width.png}
  \caption{CSS deformation half-width $aT$ as a function of recurrence
    interval $T$ for Krafla (red, $a = 5.17\times10^{-5}$~m~s$^{-1}$) and
    Asal (blue, $a = 9.47\times10^{-6}$~m~s$^{-1}$).  Vertical dashed lines
    mark estimated recurrence intervals; annotated points give the implied
    $aT$ values.  Horizontal dash-dot lines mark the total span of each
    geodetic network used in this study: 230~km for Krafla (\citealt{foulger1992post},
    1992; distance between outermost stations) and 17~km for Asal (Vigny et
    al., 2007; GPS profile extent).  In both cases $aT$ substantially exceeds
    the network aperture, explaining why inter-rifting deformation has been
    interpreted as steady plate divergence rather than a broad visco-elastic
    transient.}
  \label{fig:css_width}
\end{figure}

% ── §4.3 ─────────────────────────────────────────────────────────────────────
\subsection{Non-local diffusion versus the Fickian model}
\label{sec:disc_nonlocal}

The preference for the viscous CSS over the Foulger diffusion model
is not a numerical accident: it reflects a structural difference in the
governing equations.  \citet{foulger1992post} model post-diking deformation as
Fickian diffusion of stress, governed by $\partial u/\partial t =
\kappa\,\partial^2 u/\partial x^2$, whose Fourier signature is
$\partial\hat{u}/\partial t = -\kappa k^2\hat{u}$.  The quadratic wavenumber
dependence concentrates decay at short wavelengths and generates an erfc
displacement profile with a deformation front that propagates as
$\sqrt{\kappa t}$~and leaves the far field exponentially unperturbed on
timescales shorter than $x^2/\kappa$.  Our viscous model is governed instead
by (\ref{eq:fourier}), $\partial\hat{u}/\partial t = -a|k|\hat{u}$, where the
absolute value $|k|$ is the Fourier representation of the half-Laplacian
$(-\partial^2/\partial x^2)^{1/2}$, a spatially \emph{non-local} operator
arising from the $1/(x-x')^2$ Stokes convolution kernel (\ref{eq:stokes}).
Because the mantle transmits stress instantaneously through the incompressible
Stokes flow, every point on the elastic plate feels the rifting event from the
first instant; no diffusive front is required to reach it.  The surface
velocity is therefore a Lorentzian (\ref{eq:stepv_dimless}), with tails
decaying as $\sim 1/x^2$ rather than the exponential or Gaussian decay of the
Fickian model.

The observational consequence is directly visible in the Asal velocity
transect (Figure~\ref{fig:asal_profile}).  At stations several kilometres from
the rift axis (beyond the zone where model differences are masked by
near-field noise), the viscous CSS fits the data better than the
Foulger model ($\chi^2_r \approx 15.7$--$15.9$ vs $33$, better than a
factor-of-two improvement), because
the Lorentzian distributes deformation more broadly than the erfc profile.
The same asymmetry drives the profile misfit more clearly at Asal than at
Krafla, where the differenced-displacement observable
$\Delta u = u(\tau_f) - u(\tau_i)$ suppresses far-field signal regardless of
the spatial functional form.  The non-local character of the Stokes kernel
thus has a concrete, testable imprint on the spatial pattern of inter-rifting
deformation: networks sampling distances $x \gg aT/\pi$ are more sensitive to
the tail behaviour and therefore better placed to discriminate between the two
formulations.

The two formulations are not independent alternatives, however: the
Foulger diffusion profile is mathematically embedded within the viscous
solution. Bochner's subordination formula writes the half-Laplacian
semigroup as a weighted average, over all internal diffusion times, of the
ordinary Fickian heat semigroup \citep{metzler2000random}. Substituting the
model's own initial condition into this average reproduces the closed-form
viscous step response exactly. The Foulger model is therefore not simply
the wrong physical model: it is the correct solution evaluated at a single
diffusion timescale, rather than blended across all of them.

% ── §4.4 ─────────────────────────────────────────────────────────────────────
\subsection{Robustness under stochastic recurrence}
\label{sec:disc_stochastic}

The CSS derivation assumes strictly periodic rifting: intrusions occur at
fixed interval $T$ and deliver a fixed opening $b$ at each event.
Neither assumption holds in nature.  Magmatic rifting sequences are episodic
and irregular; the $\sim250$-year gap between the Krafla 1724--1729 and
1975--1984 crises is a single estimate from a sparse historical record, and
dike opening magnitudes vary substantially within and between sequences
\citep{thordarson2007volcanism}.  The CSS remains a robust estimator of
the long-term inter-rifting velocity field even under substantial
cycle-to-cycle variability in timing and dike opening, as we demonstrate
through stochastic ensembles.

We replaced the deterministic event sequence with stochastic ensembles in
which each past rifting event draws its recurrence interval $T_j$ and
opening $b_j$ independently from a Normal distribution with mean
$(T,\,b)$ and coefficient of variation $c_v = 0.30$, a scatter
representative of documented variability in Icelandic and Afar rift
chronologies.  For each of 500 realisations we summed the step-response
contributions from all past events and extracted the 10th--90th percentile
band across the ensemble (Figure~\ref{fig:stochastic_css}).

Four results stand out.  First, the ensemble median coincides with the
deterministic CSS in all four cases: linearity of the step-response sum and
the conditions $\mathbb{E}[T_j] = T$ and $\mathbb{E}[b_j] = b$
guarantee that the ensemble average recovers the periodic solution exactly,
regardless of the scatter level.  Second, the spread differs sharply between
the two sources of variability.  Scatter in $T$ (panel a) concentrates
uncertainty in the transition zone near $\hat{x}\sim 1$: irregular timing
shifts the velocity gradient left or right across realisations, smearing the
transition zone while the far-field plateau $v \to b/T$ remains sharp.
Scatter in $b$ (panel b) acts differently: because the recurrence
interval is fixed, events arrive at the same times in every realisation and
the profile shape is identical to the deterministic CSS in all of them;
only its overall amplitude varies.  The effect is equivalent to sliding each
realisation's curve up or down the $y$-axis by a random factor, so the band
at each distance is a fixed fraction of the local CSS value and neither the
gradient position nor the far-field plateau is affected.  Third, the combined scatter
(panel c) produces the widest band; at the deformation-width scale
$\hat{x} = 1$, the 10--90th percentile range spans roughly one-third of the
median CSS amplitude for $c_v = 0.30$, and the band narrows steadily toward
the far field as the law-of-large-numbers averaging over many past events
takes effect.  This width confirms that the CSS is a robust estimator of the
time-averaged inter-rifting velocity field even under substantial cycle-to-cycle
variability.  Fourth, imposing the plate-rate-conserving coupling
$b_j = v_{\mathrm{plate}} T_j$ (panel d) causes the far-field band to
collapse to essentially zero: because $b_j/T_j = v_{\mathrm{plate}}$
for every event in every realisation, the plate velocity is encoded exactly
regardless of timing scatter.  Residual spread in the transition zone
($|\hat{x}| \lesssim 1$) persists because stochastic timing still shifts
the position of the velocity gradient even though its far-field amplitude is
constrained.  This scenario corresponds physically to a self-regulating rift
in which crustal opening scales with the elapsed time since the previous
event, a behaviour consistent with the observation that magmatic rifts
accommodate most of the long-term plate divergence over geological timescales
\citep{thordarson2007volcanism}.

The central-limit interpretation is the following.  A geodetic snapshot of
the inter-rifting velocity field at any site within the deformation zone
$x \lesssim aT$ encodes superposed contributions from dozens of past events.
Stochastic scatter in individual $T_j$ and $b_j$ averages down over
the ensemble of contributing events, and the CSS prediction converges to
within the GPS measurement uncertainty for networks with aperture comparable
to $aT$.  Near-rift stations that respond primarily to the single most recent
event retain the full cycle-to-cycle scatter shown in panel~(a), and the CSS
approximation is less accurate there; in that regime, knowledge of the specific
rifting history is required to reduce the uncertainty.

\begin{figure}[htbp]
  \centering
  \includegraphics[width=\columnwidth]{fig/fig_stochastic_css.png}
  \caption{Robustness of the CSS velocity profile under stochastic
    recurrence ($a = T = b = 1$; phase $\hat{\tau} = 0.5$; 500
    realisations at $c_v = 0.30$).  Shaded bands show the 10th--90th
    percentile range; coloured lines are ensemble medians; the black dashed
    line is the deterministic CSS; vertical dashed lines mark
    $\hat{x} = \pm 1$.
    \textbf{(a)}~Only $T_j$ randomised.
    \textbf{(b)}~Only $b_j$ randomised.
    \textbf{(c)}~Both $T_j$ and $b_j$ randomised independently.
    \textbf{(d)}~Plate-rate-conserving coupling
    $b_j = v_{\mathrm{plate}}T_j$.
    In all four cases the ensemble median coincides with the deterministic
    CSS; see text for the physical interpretation of the band shapes.}
  \label{fig:stochastic_css}
\end{figure}

% ── §4.5 ─────────────────────────────────────────────────────────────────────
\subsection{Perspectives: the CSS as an inversion tool and its extension to
  heterogeneous rifts}
\label{sec:disc_perspectives}

The applications presented here treat $T$ and $b$ as known inputs and
fit the single free parameter $a$ to geodetic data.  The CSS framework
permits a more ambitious inversion.  Because the velocity profile
$v_{\mathrm{CSS}}(x,\tau)$ depends on three quantities (the deformation
zone half-width $aT$, the far-field plate velocity $b/T$, and the cycle
phase $\hat{\tau} = \tau/T$), a geodetic dataset that samples the near-rift
gradient and the far-field plateau simultaneously constrains all three
independently, with no locking depth or prescribed far-field forcing required.
The gradient position fixes $aT$; the long-distance plateau fixes $b/T$;
and their combination yields $a$ and $T$ individually, from which
$\eta = EH/[2a(1-\nu^2)]$ follows once $E$ and $H$ are estimated from
seismological constraints.  The stochastic analysis of
Section~\ref{sec:disc_stochastic} quantifies how cycle-to-cycle scatter in
$T_j$ and $b_j$ propagates into uncertainty on these inferred parameters,
providing a principled uncertainty budget for geodetically constrained
viscosity estimates at any magmatic rift system.

Real rifts are segmented rather than homogeneous: dike intrusions span finite
along-strike lengths, individual segments differ in recurrence history and
plate thickness, and lateral variations in heat flow modulate the effective
viscosity along the rift axis.  The 1D framework developed here is
immediately applicable to segmented geometries because its only length scale
is the local product $aT$, set by the physical properties of each segment.
Spatial variations in $H$ or $\eta$ along strike modulate $aT$ locally
without introducing new boundary conditions or requiring a finite model
domain: each segment generates its own deformation halo, and the surface
velocity field is their superposition.  The stochastic ensemble of
Section~\ref{sec:disc_stochastic} already provides the natural language for
this multi-segment setting: adjacent segments with different mean recurrence
intervals $T$ and variable event sizes $b_j$ contribute independently to
the observed velocity profile, and the CSS remains an unbiased predictor of
the long-term average field as long as $\mathbb{E}[T_j] = T$ and
$\mathbb{E}[b_j] = b$ hold locally for each segment.  Extending
the CSS to the full 2D problem (accounting for along-strike propagation
of dike intrusions and rift segment interaction) is a natural next step,
requiring the relaxation of the plane-strain approximation while retaining
the same non-local Stokes coupling at the base of the elastic plate.

% ====================================================================
\section{Conclusions}
\label{sec:conclusions}

The CSS framework resolves the kinematic gap between post-diking transients
and secular plate motion by recasting inter-rifting deformation as the
time-averaged superposition of visco-elastic step responses to a long history
of discrete dike intrusions.  Because the model satisfies far-field plate
divergence by construction and requires no locking depth, prescribed domain
width, or external velocity offset, it provides a mechanically self-consistent
alternative to backslip and dislocation models across any magmatic rift system.

At Krafla, the viscous CSS fits the post-1984 differenced displacement
record within measurement noise ($\chi^2_r = 0.72$ at $T = 200$~yr),
yielding $a = 5.17\times10^{-5}$~m~s$^{-1}$ and an implied asthenospheric
viscosity $\eta \approx 4.1\times10^{18}$~Pa~s ($H = 10$~km, $E = 40$~GPa),
consistent with post-glacial rebound estimates beneath Iceland and
reflecting the anomalously hot sub-lithospheric mantle maintained by the
Iceland plume.  Its Maxwell counterpart, with $\tau_R$ derived from the
same shared-modulus closure (\ref{eq:tau_R_shared}) rather than fit
independently, is only marginally disfavoured
($\Delta\mathrm{AIC} \approx 1.1$--$1.2$): the 8--11~yr Krafla record does
not require a Maxwell asthenosphere.  At Asal, the viscous CSS again fits
well ($\chi^2_r \approx 1.54$) and delivers $a = 9.47\times10^{-6}$~m~s$^{-1}$
and $\eta \approx 2.3\times10^{19}$~Pa~s, a factor of five higher than at
Krafla, in line with the lower heat flux of the Afar setting; here the
comparison is decisive ($\Delta\mathrm{AIC} \approx 28$--$33$ in favour of
the viscous CSS), because the 25-yr record samples enough of the
post-diking transient to rule out a spurious elastic contribution from the
derived $\tau_R$.

The non-local Stokes kernel underlying the viscous CSS produces Lorentzian
velocity tails that decay more slowly with distance than the erfc profile
of Foulger-type diffusion.  This structural difference is directly
observable in geodetic data: the viscous CSS outperforms the diffusion
model in reproducing velocities at stations several kilometres from the
rift axis, where the erfc profile underestimates the signal
($\chi^2_r \approx 15.7$--$15.9$ versus $33$ at Asal).  The Asal velocity
transect provides the clearest discrimination between the two
formulations; at Krafla, the differenced-displacement observable
$\Delta u = u(\tau_f) - u(\tau_i)$ suppresses far-field signal regardless
of the spatial functional form, so a longer or denser geodetic transect
would be required there to separate the two models statistically.

Stochastic ensembles demonstrate that the CSS is an unbiased estimator of
the long-term inter-rifting velocity field even at $c_v = 0.30$
cycle-to-cycle variability.  Scatter in recurrence interval shifts the
velocity gradient laterally, concentrating uncertainty in the transition
zone; scatter in dike opening rescales the profile amplitude uniformly,
leaving the gradient position unaffected.  Coupling the two sources through
the plate-rate constraint $b_j = v_{\mathrm{plate}} T_j$ eliminates
far-field uncertainty entirely.  In all cases the ensemble median coincides
with the deterministic CSS, confirming that stochastic recurrence introduces
spread but not bias.

These results establish the CSS as a practical tool for constraining
sub-lithospheric viscosity from geodetic inter-rifting velocity profiles
without auxiliary rheological assumptions.  Jointly inverting the gradient
position and far-field amplitude of a single velocity profile yields the
deformation zone half-width $aT$, the plate velocity $b/T$, and the
ratio $EH/\eta$ simultaneously.  Extending the framework to segmented rifts and to the full along-strike
2D geometry is the direct path toward a system-scale model of strain
accumulation at magmatic plate boundaries, retaining the
mechanically self-consistent structure derived here.

% ====================================================================
\appendix
\section{Why the step response settles at \texorpdfstring{$b/2$}{b/2}, and why \texorpdfstring{$b$}{b} remains the physical dike opening}
\label{app:b2}

Equation~(\ref{eq:step}) predicts that $u^{\text{STEP}}(x,t)\to b/2$ as
$t\to\infty$ at any fixed $x$: half the opening imposed by the initial
condition~(\ref{eq:IC}) appears to be missing. This appendix shows that the
missing half is never lost. It is a direct, necessary consequence of the
antisymmetry built into the governing equation, and the opening measured at
the intrusion site itself remains exactly $b$ for as long as the model is
evaluated. Section~\ref{app:b2_css} then extends the argument to the
Cyclo-Steady State, where the same kernel returns the full $b$, not $b/2$,
over each inter-rifting cycle, and shows why the two results are
consistent rather than contradictory.

\subsection{The initial condition splits into a stationary and an evolving part}

The Heaviside initial condition~(\ref{eq:IC}) decomposes uniquely into a
constant background and an odd perturbation about it:
\begin{equation}
u(x,0^+) = \underbrace{\frac{b}{2}}_{\text{even}}
  + \underbrace{\frac{b}{2}\,\sgn(-x)}_{\text{odd}}.
\label{eq:decomp}
\end{equation}
The governing equation~(\ref{eq:govPDE}) constrains $u$ only through its
spatial derivatives. The elastic stress~(\ref{eq:hookean}) depends on
$\partial u/\partial x$, and the Stokes traction~(\ref{eq:stokes}) on
$\partial u/\partial t$ integrated against a kernel that itself vanishes for
a spatially uniform field. A constant offset therefore carries no stress and
satisfies~(\ref{eq:govPDE}) identically at all times, so the even term in
(\ref{eq:decomp}) is an exact stationary solution, not an approximation.
Equivalently, in Fourier space~(\ref{eq:fourier_sol}) the even term is
carried entirely by the $\delta(k)$ singularity in $\hat u(k,0)$, which sits
at $k=0$ where the decay rate $a|k|$ vanishes: this is exactly the term
$I_1=b/2$ already isolated in \S\ref{sec:step}. Only the odd term, carried
by the regular part $ib/k$ of the spectrum, decays through
(\ref{eq:fourier_sol}) and produces the arctangent profile~(\ref{eq:step}).
The half that appears to vanish as $t\to\infty$ was never dynamical: it is
a rigid relabelling of the origin that the elastic--viscous system cannot
act on, because absolute displacement carries no physical meaning in an
equation that depends only on gradients of $u$.

\subsection{The opening at the source is conserved exactly}

The physically meaningful, gauge-independent quantity is not the value of
$u$ at a fixed point but the discontinuity of $u$ across the intrusion
site. Taking $x\to0^+$ and $x\to0^-$ in~(\ref{eq:step}) at any fixed $t>0$,
the argument $at/x$ diverges to $\pm\infty$ however small $t$ is, so
\begin{equation}
\lim_{x\to0^+}u^{\text{STEP}}(x,t) = \frac{b}{2}, \qquad
\lim_{x\to0^-}u^{\text{STEP}}(x,t) = -\frac{b}{2}, \qquad \forall\, t>0,
\label{eq:opening_invariant}
\end{equation}
and therefore
\begin{equation}
\lim_{x\to0^+}u^{\text{STEP}}(x,t) - \lim_{x\to0^-}u^{\text{STEP}}(x,t) = b,
\qquad \forall\, t>0.
\label{eq:opening_conserved}
\end{equation}
The opening at the source is $b$ identically, for every $t>0$, not only in
the instant after intrusion, because the odd term in (\ref{eq:decomp}) is
itself odd about $x=0$ at every later time: relaxation redistributes it in
$x$ but cannot break its antisymmetry. Nothing of the imposed slip
dissipates; the model simply reports, at $x=0$, the two one-sided limits of
a field that is continuous everywhere except at the source.

\subsection{The far-field \texorpdfstring{$b/2$}{b/2} is half of the opening having diffused outward, not the whole of it}

Equation~(\ref{eq:fourier_sol}) treats $+k$ and $-k$ identically, so the
Stokes kernel~(\ref{eq:stokes}) has no preferred direction: as the odd term
in (\ref{eq:decomp}) relaxes outward from the source, it supplies equal
strain to the $x>0$ and $x<0$ half-spaces at every instant. A given point
$x$ registers the full opening $b$ only while it remains close to the
source relative to the relaxed length scale $at$. Once $at\gg x$, that
point lies deep inside the relaxed region and has been supplied with
exactly half of the total redistributed opening, the other half having
gone to the mirror point $-x$. This is precisely what the self-similar
dependence of~(\ref{eq:step}) on $x/(at)$ encodes: for any fixed $x$, this
ratio vanishes as $t\to\infty$, so every station eventually reaches the same
$\pm b/2$ plateau regardless of how far it sits from the rift axis. The
$b/2$ asymptote is a statement about how the opening is shared between the
two flanks in the fully-relaxed limit, not about how much of it survives,
and it holds at every station, not at some special point.

This structure is not particular to the present closure. The same
$\pm(\text{slip})/2$ far-field split appears in the classical
screw-dislocation model of strain accumulation across a locked fault
\citep{savage1983dislocation}: a fault offset $s$ produces displacements of
$+s/2$ and $-s/2$ on either side, and the full offset $s$ is recovered only
as the difference between the two flanks, never as the amplitude at either
one alone. Equation~(\ref{eq:step}) inherits the same structure because
\S\ref{sec:fourier} reduces the elastic--viscous coupling to the same
non-local, parity-symmetric kernel that generates the dislocation solution.

\subsection{The CSS returns the full opening each cycle, not half of it}
\label{app:b2_css}

Section~\ref{sec:css} derives a companion result that looks, at first
sight, at odds with the $b/2$ limit above: integrating the CSS velocity
over one inter-rifting cycle at any fixed $x>0$ recovers the full opening,
$\int_0^T v_{\text{CSS}}(x,\tau)\,d\tau=b$ (\ref{eq:conservation}), not
$b/2$. The two results do not compete. They measure different quantities
on different kinds of history.

The $b/2$ plateau is the total, all-time contribution of a single, isolated
intrusion. Integrating the step velocity~(\ref{eq:stepv}) over its entire
future confirms this directly:
\begin{equation}
\int_0^{\infty}v^{\text{STEP}}(x,t)\,dt
= u^{\text{STEP}}(x,\infty)-u^{\text{STEP}}(x,0^+) = \frac{b}{2},
\qquad x>0,
\label{eq:step_flux}
\end{equation}
because $v^{\text{STEP}}\to0$ as $t\to\infty$ (\S\ref{sec:step}): once a
single event has fully relaxed, the flank comes to rest, having received
exactly $b/2$ over all of time.

The CSS never reaches that kind of rest. Because intrusions repeat every
$T$, the far-field velocity does not decay to zero but saturates on the
plate rate $v_{\text{plate}}=b/T$ (Figure~\ref{fig:css_behavior}), so
displacement at a fixed $x$ keeps advancing indefinitely instead of
settling on a plateau. What repeats from cycle to cycle is the shape of
$u_{\text{CSS}}(x,\tau)$ within $\tau\in[0,T)$, riding on an
ever-advancing baseline, and self-consistency of that repeating shape
requires the baseline to advance by exactly one dike's worth of opening,
$b$, per cycle: an advance of only $b/2$ could not close the pattern from
one intrusion to the next while also honouring (\ref{eq:step_flux}) for
every event superposed in the infinite history behind it. The full $b$
recovered in (\ref{eq:conservation}) is therefore not a doubling of the
single-event result but the signature of a state that never stops moving.
The step response gives the total distance one event carries a flank
before it comes to rest, $b/2$; the CSS gives the net distance the flank
is carried during one turn of a process that never does, $b$.

\subsection{Consequence for the fitted parameter \texorpdfstring{$b$}{b}}

Because the opening at the source is conserved exactly
(\ref{eq:opening_conserved}), the parameter $b$ recovered by fitting
(\ref{eq:step}), its cyclic sum (\ref{eq:finite_sum}), or the Maxwell and CSS
forms built on the same kernel (\S\ref{sec:maxwell_model}, \S\ref{sec:css}) to
geodetic data is the dike-opening width itself, directly comparable to
field-mapped intrusion widths without rescaling. This holds specifically for
$b$ as a model parameter, not for the displacement recorded at any single
station: a station far from the rift axis will, in the STEP long-time
limit, register only $\pm b/2$ of the opening, and comparing that
single-station amplitude against a field-mapped dike width would
underestimate it by up to a factor of two. The two quantities, the fitted
$b$ and a far-field station's own displacement, coincide only within the
near-field region where $at\lesssim|x|$.

% Note (2026-08-17): the "half-opening" label previously attached to b in
% the intro (~174), the CSS conservation statement (~520), the Maxwell CSS
% conservation statement (~769), Krafla (~1230), Asal (~1292, 1390, 1435),
% and the stochastic-recurrence section (~1626, 1638) has been removed —
% b is the dike opening everywhere, matching this appendix and the code
% (src/models.jl). The 2b appearing in the Asal discussion is unrelated to
% that label: it is the EP00-DF00 baseline-crossing factor ΔL=2u(x_0,tau)
% introduced in §sec:asal (predicted baseline spans both flanks), not a
% claim that b is half of anything.

% ====================================================================
\section*{Acknowledgements}
[PLACEHOLDER — funding agencies, field/computing support, reviewer thanks.]

% ====================================================================
\section*{Data Availability}
The geodetic data used in this study are from published sources cited in the
text. The analytical model code and inversion scripts are available at
[PLACEHOLDER — repository DOI or URL upon acceptance].

\bibliographystyle{gji}
\bibliography{refs}

\label{lastpage}
\end{document}


